\documentclass[11pt]{article}

% ---- theme variables (passed via pandoc -V) ----
\newcommand{\ThemeName}{Data Structures}
\newcommand{\ThemeKicker}{DSA Track}
\newcommand{\ThemeNumber}{02}

\usepackage{interviewstyle}
\setthemecolor{be123c}{fb7185}

% pandoc-required plumbing
\providecommand{\tightlist}{\setlength{\itemsep}{1pt}\setlength{\parskip}{0pt}}
\setlength{\emergencystretch}{3em}
\providecommand{\pandocbounded}[1]{#1}

% syntax-highlighting macros emitted by pandoc (skylighting)
\usepackage{color}
\usepackage{fancyvrb}
\newcommand{\VerbBar}{|}
\newcommand{\VERB}{\Verb[commandchars=\\\{\}]}
\DefineVerbatimEnvironment{Highlighting}{Verbatim}{commandchars=\\\{\}}
% Add ',fontsize=\small' for more characters per line
\usepackage{framed}
\definecolor{shadecolor}{RGB}{35,38,41}
\newenvironment{Shaded}{\begin{snugshade}}{\end{snugshade}}
\newcommand{\AlertTok}[1]{\textcolor[rgb]{0.58,0.85,0.30}{\textbf{\colorbox[rgb]{0.30,0.12,0.14}{#1}}}}
\newcommand{\AnnotationTok}[1]{\textcolor[rgb]{0.25,0.50,0.35}{#1}}
\newcommand{\AttributeTok}[1]{\textcolor[rgb]{0.16,0.50,0.73}{#1}}
\newcommand{\BaseNTok}[1]{\textcolor[rgb]{0.96,0.45,0.00}{#1}}
\newcommand{\BuiltInTok}[1]{\textcolor[rgb]{0.50,0.55,0.55}{#1}}
\newcommand{\CharTok}[1]{\textcolor[rgb]{0.24,0.68,0.91}{#1}}
\newcommand{\CommentTok}[1]{\textcolor[rgb]{0.48,0.49,0.49}{#1}}
\newcommand{\CommentVarTok}[1]{\textcolor[rgb]{0.50,0.55,0.55}{#1}}
\newcommand{\ConstantTok}[1]{\textcolor[rgb]{0.15,0.68,0.68}{\textbf{#1}}}
\newcommand{\ControlFlowTok}[1]{\textcolor[rgb]{0.99,0.74,0.29}{\textbf{#1}}}
\newcommand{\DataTypeTok}[1]{\textcolor[rgb]{0.16,0.50,0.73}{#1}}
\newcommand{\DecValTok}[1]{\textcolor[rgb]{0.96,0.45,0.00}{#1}}
\newcommand{\DocumentationTok}[1]{\textcolor[rgb]{0.64,0.20,0.25}{#1}}
\newcommand{\ErrorTok}[1]{\textcolor[rgb]{0.85,0.27,0.33}{\underline{#1}}}
\newcommand{\ExtensionTok}[1]{\textcolor[rgb]{0.00,0.60,1.00}{\textbf{#1}}}
\newcommand{\FloatTok}[1]{\textcolor[rgb]{0.96,0.45,0.00}{#1}}
\newcommand{\FunctionTok}[1]{\textcolor[rgb]{0.56,0.27,0.68}{#1}}
\newcommand{\ImportTok}[1]{\textcolor[rgb]{0.15,0.68,0.38}{#1}}
\newcommand{\InformationTok}[1]{\textcolor[rgb]{0.77,0.36,0.00}{#1}}
\newcommand{\KeywordTok}[1]{\textcolor[rgb]{0.81,0.81,0.76}{\textbf{#1}}}
\newcommand{\NormalTok}[1]{\textcolor[rgb]{0.81,0.81,0.76}{#1}}
\newcommand{\OperatorTok}[1]{\textcolor[rgb]{0.81,0.81,0.76}{#1}}
\newcommand{\OtherTok}[1]{\textcolor[rgb]{0.15,0.68,0.38}{#1}}
\newcommand{\PreprocessorTok}[1]{\textcolor[rgb]{0.15,0.68,0.38}{#1}}
\newcommand{\RegionMarkerTok}[1]{\textcolor[rgb]{0.16,0.50,0.73}{\colorbox[rgb]{0.08,0.19,0.26}{#1}}}
\newcommand{\SpecialCharTok}[1]{\textcolor[rgb]{0.24,0.68,0.91}{#1}}
\newcommand{\SpecialStringTok}[1]{\textcolor[rgb]{0.85,0.27,0.33}{#1}}
\newcommand{\StringTok}[1]{\textcolor[rgb]{0.96,0.31,0.31}{#1}}
\newcommand{\VariableTok}[1]{\textcolor[rgb]{0.15,0.68,0.68}{#1}}
\newcommand{\VerbatimStringTok}[1]{\textcolor[rgb]{0.85,0.27,0.33}{#1}}
\newcommand{\WarningTok}[1]{\textcolor[rgb]{0.85,0.27,0.33}{#1}}

\begin{document}

\makecoverpage

\thispagestyle{empty}
{\hypersetup{hidelinks}\tableofcontents}
\clearpage

\begin{quote}
Dense reference: mechanics, complexities, code, diagrams, and interview takeaways in one place.
\end{quote}

\section{Overview}\label{overview}

Data structures are the substrate of every algorithm. Choosing the wrong one is the single most common reason candidates fail coding interviews --- the algorithm is right, but O(n) lookups in a list where a hash map would give O(1) makes the solution too slow. This file is organized as a working reference: read top-to-bottom once, then use the complexity table and decision guide during revision.

\textbf{Mental model for selection:}

\begin{enumerate}
\def\labelenumi{\arabic{enumi}.}
\tightlist
\item
  What operations do I need? (search, insert, delete, min/max, range, prefix?)
\item
  What are the frequency/order constraints? (random access? FIFO? LIFO? sorted?)
\item
  What are the acceptable complexities?
\item
  Then pick the structure.
\end{enumerate}

\section{Master Complexity Table}\label{master-complexity-table}

\visualfigure{../../figures/dsa-02-ds/01-ds-complexity.pdf}{Pick the structure that makes your hot operation O(1)/O(log n). Hash tables win for lookups; balanced trees keep order; heaps give cheap min/max.}

All complexities are average-case unless marked. W = worst case.

\begin{longtable}[]{@{}
  >{\raggedright\arraybackslash}p{(\columnwidth - 12\tabcolsep) * \real{0.1875}}
  >{\raggedright\arraybackslash}p{(\columnwidth - 12\tabcolsep) * \real{0.1071}}
  >{\raggedright\arraybackslash}p{(\columnwidth - 12\tabcolsep) * \real{0.1071}}
  >{\raggedright\arraybackslash}p{(\columnwidth - 12\tabcolsep) * \real{0.1071}}
  >{\raggedright\arraybackslash}p{(\columnwidth - 12\tabcolsep) * \real{0.0804}}
  >{\raggedright\arraybackslash}p{(\columnwidth - 12\tabcolsep) * \real{0.0600}}
  >{\raggedright\arraybackslash}p{(\columnwidth - 12\tabcolsep) * \real{0.3571}}@{}}
\toprule\noalign{}
\rowcolor{acc!16}
\begin{minipage}[b]{\linewidth}\raggedright
Data Structure
\end{minipage} & \begin{minipage}[b]{\linewidth}\raggedright
Access
\end{minipage} & \begin{minipage}[b]{\linewidth}\raggedright
Search
\end{minipage} & \begin{minipage}[b]{\linewidth}\raggedright
Insert
\end{minipage} & \begin{minipage}[b]{\linewidth}\raggedright
Delete
\end{minipage} & \begin{minipage}[b]{\linewidth}\raggedright
Space
\end{minipage} & \begin{minipage}[b]{\linewidth}\raggedright
Notes
\end{minipage} \\
\midrule\noalign{}
\endhead
\bottomrule\noalign{}
\endlastfoot
Array (static) & O(1) & O(n) & O(n) & O(n) & O(n) & Insert/delete shifts elements \\
\rowcolor{zebra}
Dynamic Array (list) & O(1) & O(n) & O(1) amort & O(n) & O(n) & Append amortized O(1) \\
String (immutable) & O(1) & O(n) & O(n) & O(n) & O(n) & Concat is O(n); use StringBuilder \\
\rowcolor{zebra}
Hash Table & N/A & O(1) & O(1) & O(1) & O(n) & W: O(n) with all collisions \\
Singly Linked List & O(n) & O(n) & O(1) head & O(n) & O(n) & O(1) if you have the node reference \\
\rowcolor{zebra}
Doubly Linked List & O(n) & O(n) & O(1) ends & O(1)* & O(n) & *O(1) if node ref known \\
Stack & O(n) & O(n) & O(1) push & O(1) pop & O(n) & Top access O(1) \\
\rowcolor{zebra}
Queue & O(n) & O(n) & O(1) enq & O(1) deq & O(n) & Front access O(1) \\
Deque & O(1) ends & O(n) & O(1) ends & O(1) ends & O(n) & Both ends O(1) \\
\rowcolor{zebra}
BST (balanced) & O(log n) & O(log n) & O(log n) & O(log n) & O(n) & Unbalanced: O(n) worst \\
BST (unbalanced) & O(n) W & O(n) W & O(n) W & O(n) W & O(n) & Degenerate = linked list \\
\rowcolor{zebra}
AVL Tree & O(log n) & O(log n) & O(log n) & O(log n) & O(n) & Strict balance; fast lookup \\
Red-Black Tree & O(log n) & O(log n) & O(log n) & O(log n) & O(n) & Looser balance; fast insert/delete \\
\rowcolor{zebra}
Binary Heap (min/max) & O(1) top & O(n) & O(log n) & O(log n) & O(n) & Build heap: O(n) \\
Trie & O(m) & O(m) & O(m) & O(m) & O(n*m) & m = key length \\
\rowcolor{zebra}
Graph (adj list) & O(V+E) & O(V+E) & O(1) & O(E) & O(V+E) & Sparse graphs \\
Graph (adj matrix) & O(1) & O(1) edge & O(1) & O(1) & O(V\^{}2) & Dense graphs \\
\rowcolor{zebra}
Union-Find & O(\(\alpha\)(n)) & O(\(\alpha\)(n)) & O(\(\alpha\)(n)) & - & O(n) & \(\alpha\) = inverse Ackermann, \(\approx\) O(1) \\
Segment Tree & O(log n) & O(log n) & O(log n) & O(log n) & O(n) & Range queries + point updates \\
\rowcolor{zebra}
Fenwick Tree (BIT) & O(log n) & O(log n) & O(log n) & O(log n) & O(n) & Simpler than segment tree \\
LRU Cache & O(1) & O(1) & O(1) & O(1) & O(n) & Hash map + doubly linked list \\
\rowcolor{zebra}
Bloom Filter & N/A & O(k) & O(k) & N/A & O(m) & k=hash funcs, probabilistic \\
\end{longtable}

\section{The Data Structures}\label{the-data-structures}

\subsection{Arrays \& Dynamic Arrays}\label{arrays--dynamic-arrays}

\textbf{What it is:} Contiguous block of memory where each element sits at a fixed offset from the base address. Random access is O(1) because address = base + index * element\_size.

\textbf{Internal mechanics:}

\setcodelang{}

\begin{verbatim}
Static Array (size=6):
Index:  [0]  [1]  [2]  [3]  [4]  [5]
Value:  [10] [20] [30] [40] [50] [60]
         ^                         ^
         base addr                 base + 5*sizeof(int)

Dynamic Array (Python list) — amortized resize:
size=4, capacity=4:   [1][2][3][4][ ][ ][ ][ ]  ← capacity doubles on overflow
Append 5 → capacity=8:[1][2][3][4][5][ ][ ][ ]
Append cost analysis (n appends):
  - n-1 cheap appends: O(1) each
  - ~log(n) doublings: cost 1+2+4+...+n = 2n copies total
  - Amortized per append: 2n / n = O(1) ✓
\end{verbatim}

\textbf{Complexity:}

\begin{longtable}[]{@{}
  >{\raggedright\arraybackslash}p{(\columnwidth - 6\tabcolsep) * \real{0.2948}}
  >{\raggedright\arraybackslash}p{(\columnwidth - 6\tabcolsep) * \real{0.2260}}
  >{\raggedright\arraybackslash}p{(\columnwidth - 6\tabcolsep) * \real{0.1966}}
  >{\raggedright\arraybackslash}p{(\columnwidth - 6\tabcolsep) * \real{0.2826}}@{}}
\toprule\noalign{}
\rowcolor{acc!16}
\begin{minipage}[b]{\linewidth}\raggedright
Op
\end{minipage} & \begin{minipage}[b]{\linewidth}\raggedright
Best
\end{minipage} & \begin{minipage}[b]{\linewidth}\raggedright
Avg
\end{minipage} & \begin{minipage}[b]{\linewidth}\raggedright
Worst
\end{minipage} \\
\midrule\noalign{}
\endhead
\bottomrule\noalign{}
\endlastfoot
Access & O(1) & O(1) & O(1) \\
\rowcolor{zebra}
Search & O(1) & O(n) & O(n) \\
Append & O(1) & O(1) & O(n) \\
\rowcolor{zebra}
Insert & O(1) & O(n) & O(n) \\
Delete & O(1) & O(n) & O(n) \\
\end{longtable}

\textbf{When to use:}

\begin{itemize}
\tightlist
\item
  Random access by index is needed
\item
  Data is mostly read; writes are appends
\item
  Cache locality matters (iteration speed)
\end{itemize}

\textbf{When NOT to use:}

\begin{itemize}
\tightlist
\item
  Frequent mid-list insertions/deletions (use linked list or deque)
\item
  Need fast search (use hash map or BST)
\end{itemize}

\textbf{Key Python patterns:}

\setcodelang{Python}

\begin{Shaded}
\begin{Highlighting}[]
\CommentTok{\# Two{-}pointer on sorted array}
\NormalTok{left, right }\OperatorTok{=} \DecValTok{0}\NormalTok{, }\BuiltInTok{len}\NormalTok{(arr) }\OperatorTok{{-}} \DecValTok{1}
\ControlFlowTok{while}\NormalTok{ left }\OperatorTok{\textless{}}\NormalTok{ right:}
\NormalTok{    s }\OperatorTok{=}\NormalTok{ arr[left] }\OperatorTok{+}\NormalTok{ arr[right]}
    \ControlFlowTok{if}\NormalTok{ s }\OperatorTok{==}\NormalTok{ target: }\ControlFlowTok{return}\NormalTok{ [left, right]}
    \ControlFlowTok{elif}\NormalTok{ s }\OperatorTok{\textless{}}\NormalTok{ target: left }\OperatorTok{+=} \DecValTok{1}
    \ControlFlowTok{else}\NormalTok{: right }\OperatorTok{{-}=} \DecValTok{1}

\CommentTok{\# Sliding window}
\NormalTok{window\_sum }\OperatorTok{=} \BuiltInTok{sum}\NormalTok{(arr[:k])}
\NormalTok{max\_sum }\OperatorTok{=}\NormalTok{ window\_sum}
\ControlFlowTok{for}\NormalTok{ i }\KeywordTok{in} \BuiltInTok{range}\NormalTok{(k, }\BuiltInTok{len}\NormalTok{(arr)):}
\NormalTok{    window\_sum }\OperatorTok{+=}\NormalTok{ arr[i] }\OperatorTok{{-}}\NormalTok{ arr[i }\OperatorTok{{-}}\NormalTok{ k]}
\NormalTok{    max\_sum }\OperatorTok{=} \BuiltInTok{max}\NormalTok{(max\_sum, window\_sum)}

\CommentTok{\# In{-}place reversal}
\NormalTok{arr[i], arr[j] }\OperatorTok{=}\NormalTok{ arr[j], arr[i]}
\end{Highlighting}
\end{Shaded}

\textbf{Interview takeaway:} \textbf{Arrays are the default. When the problem says "return sorted" or "in-place", think array + two pointers. When you need to insert/delete mid-sequence often, question whether an array is right.}

\textbf{Real interview use-case:} Two Sum (hash map), Sliding Window Maximum, Kadane\textquotesingle s (max subarray), rotate array in-place, merge sorted arrays.

\subsection{Strings}\label{strings}

\textbf{What it is:} Sequence of characters. In Python and Java, strings are \textbf{immutable} --- every concatenation creates a new object. In C++, \texttt{std::string} is mutable but each \texttt{+=} may reallocate.

\textbf{Internal mechanics:}

\setcodelang{}

\begin{verbatim}
Python string "hello":
  h | e | l | l | o
  Each character is a Unicode code point (Python 3: variable width internally)
  id("hello" + "x") != id("hello")  ← new object every concat

String concatenation cost in a loop (naïve):
  for i in range(n): result += s[i]
  Iteration 1: copy 1 char    → O(1)
  Iteration 2: copy 2 chars   → O(2)
  ...
  Iteration n: copy n chars   → O(n)
  Total: O(1+2+...+n) = O(n²)  ← DISASTER for large n

String Builder pattern:
  parts = []
  for c in chars: parts.append(c)   ← O(1) amortized each
  result = "".join(parts)            ← O(n) one pass
  Total: O(n)  ✓
\end{verbatim}

\textbf{Complexity:}

\begin{longtable}[]{@{}
  >{\raggedright\arraybackslash}p{(\columnwidth - 4\tabcolsep) * \real{0.2474}}
  >{\raggedright\arraybackslash}p{(\columnwidth - 4\tabcolsep) * \real{0.2371}}
  >{\raggedright\arraybackslash}p{(\columnwidth - 4\tabcolsep) * \real{0.5155}}@{}}
\toprule\noalign{}
\rowcolor{acc!16}
\begin{minipage}[b]{\linewidth}\raggedright
Op
\end{minipage} & \begin{minipage}[b]{\linewidth}\raggedright
Complexity
\end{minipage} & \begin{minipage}[b]{\linewidth}\raggedright
Note
\end{minipage} \\
\midrule\noalign{}
\endhead
\bottomrule\noalign{}
\endlastfoot
Index access & O(1) & \\
\rowcolor{zebra}
Slice s{[}i:j{]} & O(j-i) & Creates new string \\
Concatenate & O(n+m) & Both strings copied \\
\rowcolor{zebra}
Find/search & O(n*m) & Naive; KMP/Z gives O(n+m) \\
Comparison & O(n) & Character by character \\
\rowcolor{zebra}
Join & O(n) & One pass \\
\end{longtable}

\textbf{Key Python patterns:}

\setcodelang{Python}

\begin{Shaded}
\begin{Highlighting}[]
\CommentTok{\# Always join, never += in loop}
\NormalTok{result }\OperatorTok{=} \StringTok{""}\NormalTok{.join(list\_of\_chars)}

\CommentTok{\# Reverse a string}
\NormalTok{rev }\OperatorTok{=}\NormalTok{ s[::}\OperatorTok{{-}}\DecValTok{1}\NormalTok{]}

\CommentTok{\# Check palindrome}
\KeywordTok{def}\NormalTok{ is\_palindrome(s):}
\NormalTok{    s }\OperatorTok{=}\NormalTok{ s.lower()}
\NormalTok{    l, r }\OperatorTok{=} \DecValTok{0}\NormalTok{, }\BuiltInTok{len}\NormalTok{(s) }\OperatorTok{{-}} \DecValTok{1}
    \ControlFlowTok{while}\NormalTok{ l }\OperatorTok{\textless{}}\NormalTok{ r:}
        \ControlFlowTok{while}\NormalTok{ l }\OperatorTok{\textless{}}\NormalTok{ r }\KeywordTok{and} \KeywordTok{not}\NormalTok{ s[l].isalnum(): l }\OperatorTok{+=} \DecValTok{1}
        \ControlFlowTok{while}\NormalTok{ l }\OperatorTok{\textless{}}\NormalTok{ r }\KeywordTok{and} \KeywordTok{not}\NormalTok{ s[r].isalnum(): r }\OperatorTok{{-}=} \DecValTok{1}
        \ControlFlowTok{if}\NormalTok{ s[l] }\OperatorTok{!=}\NormalTok{ s[r]: }\ControlFlowTok{return} \VariableTok{False}
\NormalTok{        l }\OperatorTok{+=} \DecValTok{1}\OperatorTok{;}\NormalTok{ r }\OperatorTok{{-}=} \DecValTok{1}
    \ControlFlowTok{return} \VariableTok{True}

\CommentTok{\# Anagram check}
\ImportTok{from}\NormalTok{ collections }\ImportTok{import}\NormalTok{ Counter}
\NormalTok{Counter(s) }\OperatorTok{==}\NormalTok{ Counter(t)}

\CommentTok{\# Sliding window on string (distinct chars)}
\ImportTok{from}\NormalTok{ collections }\ImportTok{import}\NormalTok{ defaultdict}
\NormalTok{freq }\OperatorTok{=}\NormalTok{ defaultdict(}\BuiltInTok{int}\NormalTok{)}
\NormalTok{l }\OperatorTok{=} \DecValTok{0}
\ControlFlowTok{for}\NormalTok{ r, c }\KeywordTok{in} \BuiltInTok{enumerate}\NormalTok{(s):}
\NormalTok{    freq[c] }\OperatorTok{+=} \DecValTok{1}
    \ControlFlowTok{while} \BuiltInTok{len}\NormalTok{(freq) }\OperatorTok{\textgreater{}}\NormalTok{ k:}
\NormalTok{        freq[s[l]] }\OperatorTok{{-}=} \DecValTok{1}
        \ControlFlowTok{if}\NormalTok{ freq[s[l]] }\OperatorTok{==} \DecValTok{0}\NormalTok{: }\KeywordTok{del}\NormalTok{ freq[s[l]]}
\NormalTok{        l }\OperatorTok{+=} \DecValTok{1}
\end{Highlighting}
\end{Shaded}

\textbf{Interview takeaway:} \textbf{Never concatenate strings in a loop. Use \texttt{"".join(parts)}. For pattern matching, think sliding window + frequency map. Anagram/permutation problems \(\rightarrow\) character frequency counter.}

\textbf{Real interview use-case:} Longest substring without repeating characters, group anagrams, valid palindrome, minimum window substring, encode/decode strings.

\subsection{Hash Tables}\label{hash-tables}

\textbf{What it is:} Maps keys to values in O(1) average by computing \texttt{hash(key)\ \%\ capacity} to find a bucket, then storing the key-value pair there.

\textbf{Internal mechanics:}

\setcodelang{}

\begin{verbatim}
Hash Table with Chaining (load factor = 0.75):

  Buckets:
  [0] → None
  [1] → ("apple", 5) → ("grape", 2) → None   ← collision! same bucket
  [2] → ("banana", 3) → None
  [3] → None
  [4] → ("cherry", 7) → None
  ...

  hash("apple")  % 8 = 1
  hash("grape")  % 8 = 1   ← collision → chain extends
  hash("banana") % 8 = 2

Open Addressing (linear probing):
  [0]: empty
  [1]: ("apple", 5)
  [2]: ("grape", 2)   ← collision at 1, probe to 2
  [3]: ("banana", 3)
  ...

Load Factor (α) = n / capacity
  α < 0.7: few collisions, O(1) expected
  α > 0.9: many collisions → O(n) worst
  → resize (rehash all) when α threshold exceeded
  Python dict resizes at 2/3 full → rehash to 2x capacity

Good Hash Function Properties:
  1. Deterministic (same key → same hash)
  2. Uniform distribution (minimize collisions)
  3. Fast to compute
  4. Avalanche effect (small key change → big hash change)
\end{verbatim}

\textbf{Collision resolution comparison:}

\begin{longtable}[]{@{}
  >{\raggedright\arraybackslash}p{(\columnwidth - 4\tabcolsep) * \real{0.2083}}
  >{\raggedright\arraybackslash}p{(\columnwidth - 4\tabcolsep) * \real{0.3611}}
  >{\raggedright\arraybackslash}p{(\columnwidth - 4\tabcolsep) * \real{0.4306}}@{}}
\toprule\noalign{}
\rowcolor{acc!16}
\begin{minipage}[b]{\linewidth}\raggedright
Strategy
\end{minipage} & \begin{minipage}[b]{\linewidth}\raggedright
Pro
\end{minipage} & \begin{minipage}[b]{\linewidth}\raggedright
Con
\end{minipage} \\
\midrule\noalign{}
\endhead
\bottomrule\noalign{}
\endlastfoot
Chaining & Simple, works at high load & Extra memory for pointers \\
\rowcolor{zebra}
Linear probing & Cache-friendly & Clustering degrades performance \\
Quadratic probe & Less clustering & May not probe all slots \\
\rowcolor{zebra}
Double hashing & Best distribution & Two hash functions needed \\
\end{longtable}

\textbf{Complexity:}

\begin{longtable}[]{@{}
  >{\raggedright\arraybackslash}p{(\columnwidth - 4\tabcolsep) * \real{0.3030}}
  >{\raggedright\arraybackslash}p{(\columnwidth - 4\tabcolsep) * \real{0.4066}}
  >{\raggedright\arraybackslash}p{(\columnwidth - 4\tabcolsep) * \real{0.2904}}@{}}
\toprule\noalign{}
\rowcolor{acc!16}
\begin{minipage}[b]{\linewidth}\raggedright
Op
\end{minipage} & \begin{minipage}[b]{\linewidth}\raggedright
Average
\end{minipage} & \begin{minipage}[b]{\linewidth}\raggedright
Worst
\end{minipage} \\
\midrule\noalign{}
\endhead
\bottomrule\noalign{}
\endlastfoot
Search & O(1) & O(n) \\
\rowcolor{zebra}
Insert & O(1) & O(n) \\
Delete & O(1) & O(n) \\
\end{longtable}

\textbf{When to use:}

\begin{itemize}
\tightlist
\item
  O(1) lookup by key
\item
  Counting frequencies
\item
  Caching / memoization
\item
  Deduplication
\end{itemize}

\textbf{When NOT to use:}

\begin{itemize}
\tightlist
\item
  Need sorted order (use BST/sorted array)
\item
  Keys not hashable (use tree map)
\item
  Memory is very tight
\end{itemize}

\textbf{Key Python patterns:}

\setcodelang{Python}

\begin{Shaded}
\begin{Highlighting}[]
\ImportTok{from}\NormalTok{ collections }\ImportTok{import}\NormalTok{ defaultdict, Counter}

\CommentTok{\# Frequency count}
\NormalTok{freq }\OperatorTok{=}\NormalTok{ Counter(}\StringTok{"anagram"}\NormalTok{)            }\CommentTok{\# Counter(\{\textquotesingle{}a\textquotesingle{}: 3, \textquotesingle{}n\textquotesingle{}: 1, ...\})}

\CommentTok{\# Default dict avoids KeyError}
\NormalTok{graph }\OperatorTok{=}\NormalTok{ defaultdict(}\BuiltInTok{list}\NormalTok{)}
\NormalTok{graph[}\StringTok{\textquotesingle{}A\textquotesingle{}}\NormalTok{].append(}\StringTok{\textquotesingle{}B\textquotesingle{}}\NormalTok{)}

\CommentTok{\# Two Sum O(n)}
\KeywordTok{def}\NormalTok{ two\_sum(nums, target):}
\NormalTok{    seen }\OperatorTok{=}\NormalTok{ \{\}}
    \ControlFlowTok{for}\NormalTok{ i, n }\KeywordTok{in} \BuiltInTok{enumerate}\NormalTok{(nums):}
\NormalTok{        complement }\OperatorTok{=}\NormalTok{ target }\OperatorTok{{-}}\NormalTok{ n}
        \ControlFlowTok{if}\NormalTok{ complement }\KeywordTok{in}\NormalTok{ seen:}
            \ControlFlowTok{return}\NormalTok{ [seen[complement], i]}
\NormalTok{        seen[n] }\OperatorTok{=}\NormalTok{ i}

\CommentTok{\# Grouping anagrams}
\ImportTok{from}\NormalTok{ collections }\ImportTok{import}\NormalTok{ defaultdict}
\KeywordTok{def}\NormalTok{ group\_anagrams(strs):}
\NormalTok{    groups }\OperatorTok{=}\NormalTok{ defaultdict(}\BuiltInTok{list}\NormalTok{)}
    \ControlFlowTok{for}\NormalTok{ s }\KeywordTok{in}\NormalTok{ strs:}
\NormalTok{        groups[}\BuiltInTok{tuple}\NormalTok{(}\BuiltInTok{sorted}\NormalTok{(s))].append(s)}
    \ControlFlowTok{return} \BuiltInTok{list}\NormalTok{(groups.values())}
\end{Highlighting}
\end{Shaded}

\textbf{Interview takeaway:} \textbf{Hash map is your first tool when you need O(1) lookup. "Two pointer O(n\(^2\))? Can we trade space for time?" \(\rightarrow\) hash map. Watch out: dict keys must be hashable (no lists as keys; use tuples).}

\textbf{Real interview use-case:} Two Sum, subarray sum equals k, longest consecutive sequence, top K frequent elements, valid sudoku.

\subsection{Linked Lists}\label{linked-lists}

\textbf{What it is:} Sequence of nodes where each node holds a value and a pointer to the next (singly) or both next and previous (doubly) nodes. Non-contiguous memory.

\textbf{Internal mechanics:}

\setcodelang{}

\begin{verbatim}
Singly Linked List:
  head
   ↓
  [1|•]→[2|•]→[3|•]→[4|•]→None
   node   next pointer

Doubly Linked List:
  head                              tail
   ↓                                ↓
None←[•|1|•]↔[•|2|•]↔[•|3|•]↔[•|4|•]→None
     prev val next

Insert at head (O(1)):
  new_node.next = head
  head = new_node

Insert at tail (O(1) with tail pointer):
  tail.next = new_node
  tail = new_node

Delete node (given reference, doubly linked, O(1)):
  node.prev.next = node.next
  node.next.prev = node.prev

Delete node (singly linked, given reference only, O(n)):
  Must traverse from head to find predecessor
  → trick: copy next node's value into current, delete next
\end{verbatim}

\textbf{Complexity:}

\begin{longtable}[]{@{}
  >{\raggedright\arraybackslash}p{(\columnwidth - 4\tabcolsep) * \real{0.4360}}
  >{\raggedright\arraybackslash}p{(\columnwidth - 4\tabcolsep) * \real{0.2820}}
  >{\raggedright\arraybackslash}p{(\columnwidth - 4\tabcolsep) * \real{0.2820}}@{}}
\toprule\noalign{}
\rowcolor{acc!16}
\begin{minipage}[b]{\linewidth}\raggedright
Op
\end{minipage} & \begin{minipage}[b]{\linewidth}\raggedright
Singly LL
\end{minipage} & \begin{minipage}[b]{\linewidth}\raggedright
Doubly LL
\end{minipage} \\
\midrule\noalign{}
\endhead
\bottomrule\noalign{}
\endlastfoot
Access by index & O(n) & O(n) \\
\rowcolor{zebra}
Search & O(n) & O(n) \\
Insert at head & O(1) & O(1) \\
\rowcolor{zebra}
Insert at tail & O(n)/O(1) & O(1) \\
Insert at middle & O(n) & O(n) \\
\rowcolor{zebra}
Delete at head & O(1) & O(1) \\
Delete by ref & O(n) & O(1) \\
\end{longtable}

\textbf{When to use:}

\begin{itemize}
\tightlist
\item
  Frequent insertions/deletions at head or known positions
\item
  Implementing stacks, queues, LRU cache
\item
  When you don\textquotesingle t need random access
\end{itemize}

\textbf{When NOT to use:}

\begin{itemize}
\tightlist
\item
  Random access needed (use array)
\item
  Memory is tight (each node has pointer overhead)
\item
  Cache performance critical (non-contiguous = cache misses)
\end{itemize}

\textbf{Key Python patterns:}

\setcodelang{Python}

\begin{Shaded}
\begin{Highlighting}[]
\KeywordTok{class}\NormalTok{ ListNode:}
    \KeywordTok{def} \FunctionTok{\_\_init\_\_}\NormalTok{(}\VariableTok{self}\NormalTok{, val}\OperatorTok{=}\DecValTok{0}\NormalTok{, }\BuiltInTok{next}\OperatorTok{=}\VariableTok{None}\NormalTok{):}
        \VariableTok{self}\NormalTok{.val }\OperatorTok{=}\NormalTok{ val}
        \VariableTok{self}\NormalTok{.}\BuiltInTok{next} \OperatorTok{=} \BuiltInTok{next}

\CommentTok{\# Reverse a linked list (iterative)}
\KeywordTok{def}\NormalTok{ reverse(head):}
\NormalTok{    prev }\OperatorTok{=} \VariableTok{None}
\NormalTok{    curr }\OperatorTok{=}\NormalTok{ head}
    \ControlFlowTok{while}\NormalTok{ curr:}
\NormalTok{        nxt }\OperatorTok{=}\NormalTok{ curr.}\BuiltInTok{next}
\NormalTok{        curr.}\BuiltInTok{next} \OperatorTok{=}\NormalTok{ prev}
\NormalTok{        prev }\OperatorTok{=}\NormalTok{ curr}
\NormalTok{        curr }\OperatorTok{=}\NormalTok{ nxt}
    \ControlFlowTok{return}\NormalTok{ prev}

\CommentTok{\# Floyd\textquotesingle{}s cycle detection}
\KeywordTok{def}\NormalTok{ has\_cycle(head):}
\NormalTok{    slow }\OperatorTok{=}\NormalTok{ fast }\OperatorTok{=}\NormalTok{ head}
    \ControlFlowTok{while}\NormalTok{ fast }\KeywordTok{and}\NormalTok{ fast.}\BuiltInTok{next}\NormalTok{:}
\NormalTok{        slow }\OperatorTok{=}\NormalTok{ slow.}\BuiltInTok{next}
\NormalTok{        fast }\OperatorTok{=}\NormalTok{ fast.}\BuiltInTok{next}\NormalTok{.}\BuiltInTok{next}
        \ControlFlowTok{if}\NormalTok{ slow }\KeywordTok{is}\NormalTok{ fast:}
            \ControlFlowTok{return} \VariableTok{True}
    \ControlFlowTok{return} \VariableTok{False}

\CommentTok{\# Find middle (slow/fast pointers)}
\KeywordTok{def}\NormalTok{ find\_middle(head):}
\NormalTok{    slow }\OperatorTok{=}\NormalTok{ fast }\OperatorTok{=}\NormalTok{ head}
    \ControlFlowTok{while}\NormalTok{ fast }\KeywordTok{and}\NormalTok{ fast.}\BuiltInTok{next}\NormalTok{:}
\NormalTok{        slow }\OperatorTok{=}\NormalTok{ slow.}\BuiltInTok{next}
\NormalTok{        fast }\OperatorTok{=}\NormalTok{ fast.}\BuiltInTok{next}\NormalTok{.}\BuiltInTok{next}
    \ControlFlowTok{return}\NormalTok{ slow}

\CommentTok{\# Dummy node pattern (simplifies edge cases)}
\KeywordTok{def}\NormalTok{ delete\_nth\_from\_end(head, n):}
\NormalTok{    dummy }\OperatorTok{=}\NormalTok{ ListNode(}\DecValTok{0}\NormalTok{, head)}
\NormalTok{    fast }\OperatorTok{=}\NormalTok{ slow }\OperatorTok{=}\NormalTok{ dummy}
    \ControlFlowTok{for}\NormalTok{ \_ }\KeywordTok{in} \BuiltInTok{range}\NormalTok{(n }\OperatorTok{+} \DecValTok{1}\NormalTok{):}
\NormalTok{        fast }\OperatorTok{=}\NormalTok{ fast.}\BuiltInTok{next}
    \ControlFlowTok{while}\NormalTok{ fast:}
\NormalTok{        fast }\OperatorTok{=}\NormalTok{ fast.}\BuiltInTok{next}
\NormalTok{        slow }\OperatorTok{=}\NormalTok{ slow.}\BuiltInTok{next}
\NormalTok{    slow.}\BuiltInTok{next} \OperatorTok{=}\NormalTok{ slow.}\BuiltInTok{next}\NormalTok{.}\BuiltInTok{next}
    \ControlFlowTok{return}\NormalTok{ dummy.}\BuiltInTok{next}
\end{Highlighting}
\end{Shaded}

\textbf{Interview takeaway:} \textbf{Always use a dummy head node to avoid null checks at head. Slow/fast pointer solves cycle detection, middle finding, and nth-from-end in one pass. Linked list problems are usually just pointer manipulation --- draw it out.}

\textbf{Real interview use-case:} Reverse linked list, merge two sorted lists, detect cycle, reorder list, LRU cache.

\subsection{Stacks \& Queues}\label{stacks--queues}

\textbf{What it is:}

\begin{itemize}
\tightlist
\item
  \textbf{Stack}: LIFO (Last In, First Out). Like a stack of plates.
\item
  \textbf{Queue}: FIFO (First In, First Out). Like a line at a store.
\item
  \textbf{Deque}: Double-ended queue; insert/delete at both ends O(1).
\item
  \textbf{Monotonic stack/queue}: Stack/queue that maintains monotone order for range queries.
\item
  \textbf{Circular buffer}: Fixed-size ring buffer for queue; avoids shifting.
\end{itemize}

\textbf{Internal mechanics:}

\setcodelang{}

\begin{verbatim}
Stack (LIFO):
  Push 1,2,3:  [1]→[2]→[3]   ← top is 3
  Pop:         [1]→[2]        ← returns 3

Queue (FIFO) using circular buffer:
  Capacity=5, head=0, tail=0
  Enqueue 1,2,3:
    [1][2][3][ ][ ]   head=0, tail=3
  Dequeue:
    [ ][2][3][ ][ ]   head=1, tail=3
  Enqueue 4,5:
    [ ][2][3][4][5]   head=1, tail=0 (wrapped!)
  → indices mod capacity: (tail+1) % cap

Monotonic Stack (next greater element):
  arr = [2, 1, 5, 3, 4]
  Stack tracks indices of elements waiting for their "next greater"
  i=0: push 0       stack=[0]         (arr[0]=2)
  i=1: 1<2, push 1  stack=[0,1]       (arr[1]=1)
  i=2: 5>1, pop 1→ans[1]=5, 5>2 pop 0→ans[0]=5, push 2
                    stack=[2]
  i=3: 3<5, push 3  stack=[2,3]
  i=4: 4>3, pop 3→ans[3]=4, 4<5 push 4
                    stack=[2,4]
  Remaining: no greater element → -1
  ans = [5, 5, -1, 4, -1]

Deque (collections.deque in Python):
  appendleft / popleft  → O(1)
  append / pop          → O(1)
\end{verbatim}

\textbf{Complexity:}

\begin{longtable}[]{@{}
  >{\raggedright\arraybackslash}p{(\columnwidth - 8\tabcolsep) * \real{0.3168}}
  >{\raggedright\arraybackslash}p{(\columnwidth - 8\tabcolsep) * \real{0.1139}}
  >{\raggedright\arraybackslash}p{(\columnwidth - 8\tabcolsep) * \real{0.1139}}
  >{\raggedright\arraybackslash}p{(\columnwidth - 8\tabcolsep) * \real{0.1139}}
  >{\raggedright\arraybackslash}p{(\columnwidth - 8\tabcolsep) * \real{0.3416}}@{}}
\toprule\noalign{}
\rowcolor{acc!16}
\begin{minipage}[b]{\linewidth}\raggedright
Op
\end{minipage} & \begin{minipage}[b]{\linewidth}\raggedright
Stack
\end{minipage} & \begin{minipage}[b]{\linewidth}\raggedright
Queue
\end{minipage} & \begin{minipage}[b]{\linewidth}\raggedright
Deque
\end{minipage} & \begin{minipage}[b]{\linewidth}\raggedright
Monotonic Stack
\end{minipage} \\
\midrule\noalign{}
\endhead
\bottomrule\noalign{}
\endlastfoot
Push/Enqueue & O(1) & O(1) & O(1) & O(1) amort \\
\rowcolor{zebra}
Pop/Dequeue & O(1) & O(1) & O(1) & O(1) amort \\
Peek (top/front) & O(1) & O(1) & O(1) & O(1) \\
\rowcolor{zebra}
Search & O(n) & O(n) & O(n) & O(n) \\
\end{longtable}

\textbf{When to use:}

\begin{itemize}
\tightlist
\item
  Stack: DFS, expression parsing, undo operations, backtracking
\item
  Queue: BFS, task scheduling, sliding window with deque
\item
  Monotonic stack: next greater/smaller element, largest rectangle in histogram
\item
  Deque: sliding window maximum, palindrome check
\end{itemize}

\textbf{Key Python patterns:}

\setcodelang{Python}

\begin{Shaded}
\begin{Highlighting}[]
\ImportTok{from}\NormalTok{ collections }\ImportTok{import}\NormalTok{ deque}

\CommentTok{\# Stack (use list, pop from end)}
\NormalTok{stack }\OperatorTok{=}\NormalTok{ []}
\NormalTok{stack.append(x)    }\CommentTok{\# push}
\NormalTok{stack.pop()        }\CommentTok{\# pop O(1)}
\NormalTok{stack[}\OperatorTok{{-}}\DecValTok{1}\NormalTok{]          }\CommentTok{\# peek}

\CommentTok{\# Queue (use deque for O(1) dequeue)}
\NormalTok{q }\OperatorTok{=}\NormalTok{ deque()}
\NormalTok{q.append(x)        }\CommentTok{\# enqueue}
\NormalTok{q.popleft()        }\CommentTok{\# dequeue O(1)}

\CommentTok{\# BFS template}
\KeywordTok{def}\NormalTok{ bfs(root):}
\NormalTok{    q }\OperatorTok{=}\NormalTok{ deque([root])}
    \ControlFlowTok{while}\NormalTok{ q:}
\NormalTok{        node }\OperatorTok{=}\NormalTok{ q.popleft()}
        \ControlFlowTok{for}\NormalTok{ child }\KeywordTok{in}\NormalTok{ node.children:}
\NormalTok{            q.append(child)}

\CommentTok{\# Monotonic stack — next greater element}
\KeywordTok{def}\NormalTok{ next\_greater(nums):}
\NormalTok{    n }\OperatorTok{=} \BuiltInTok{len}\NormalTok{(nums)}
\NormalTok{    ans }\OperatorTok{=}\NormalTok{ [}\OperatorTok{{-}}\DecValTok{1}\NormalTok{] }\OperatorTok{*}\NormalTok{ n}
\NormalTok{    stack }\OperatorTok{=}\NormalTok{ []  }\CommentTok{\# indices}
    \ControlFlowTok{for}\NormalTok{ i, num }\KeywordTok{in} \BuiltInTok{enumerate}\NormalTok{(nums):}
        \ControlFlowTok{while}\NormalTok{ stack }\KeywordTok{and}\NormalTok{ nums[stack[}\OperatorTok{{-}}\DecValTok{1}\NormalTok{]] }\OperatorTok{\textless{}}\NormalTok{ num:}
\NormalTok{            idx }\OperatorTok{=}\NormalTok{ stack.pop()}
\NormalTok{            ans[idx] }\OperatorTok{=}\NormalTok{ num}
\NormalTok{        stack.append(i)}
    \ControlFlowTok{return}\NormalTok{ ans}

\CommentTok{\# Sliding window maximum (monotonic deque)}
\KeywordTok{def}\NormalTok{ sliding\_max(nums, k):}
\NormalTok{    dq }\OperatorTok{=}\NormalTok{ deque()   }\CommentTok{\# indices, decreasing values}
\NormalTok{    res }\OperatorTok{=}\NormalTok{ []}
    \ControlFlowTok{for}\NormalTok{ i, num }\KeywordTok{in} \BuiltInTok{enumerate}\NormalTok{(nums):}
        \ControlFlowTok{while}\NormalTok{ dq }\KeywordTok{and}\NormalTok{ nums[dq[}\OperatorTok{{-}}\DecValTok{1}\NormalTok{]] }\OperatorTok{\textless{}}\NormalTok{ num:}
\NormalTok{            dq.pop()}
\NormalTok{        dq.append(i)}
        \ControlFlowTok{if}\NormalTok{ dq[}\DecValTok{0}\NormalTok{] }\OperatorTok{==}\NormalTok{ i }\OperatorTok{{-}}\NormalTok{ k:  }\CommentTok{\# out of window}
\NormalTok{            dq.popleft()}
        \ControlFlowTok{if}\NormalTok{ i }\OperatorTok{\textgreater{}=}\NormalTok{ k }\OperatorTok{{-}} \DecValTok{1}\NormalTok{:}
\NormalTok{            res.append(nums[dq[}\DecValTok{0}\NormalTok{]])}
    \ControlFlowTok{return}\NormalTok{ res}
\end{Highlighting}
\end{Shaded}

\textbf{Interview takeaway:} \textbf{Monotonic stack is the pattern for "next greater/smaller element" and "largest rectangle" problems. BFS always uses a queue. For sliding window max/min, use a monotonic deque. Never use \texttt{list.pop(0)} for queues --- it\textquotesingle s O(n); use \texttt{deque.popleft()}.}

\textbf{Real interview use-case:} Valid parentheses, min stack, daily temperatures, largest rectangle in histogram, sliding window maximum.

\subsection{Trees}\label{trees}

\textbf{What it is:} Hierarchical, acyclic connected graph. Root at top, leaves at bottom.

\textbf{Internal mechanics:}

\setcodelang{}

\begin{verbatim}
Binary Tree (each node: up to 2 children):
           1
          / \
         2   3
        / \   \
       4   5   6

BST (Binary Search Tree): left < root < right
           8
          / \
         3   10
        / \    \
       1   6   14
          / \  /
         4   7 13

BST invariant: for any node N:
  - All nodes in left subtree < N.val
  - All nodes in right subtree > N.val
  → In-order traversal gives sorted sequence!

BST Search:
  Find 6: start at 8 → 6<8 go left → 3 → 6>3 go right → 6 ✓
  O(h) where h = height. Balanced: h=log n. Skewed: h=n.

Balanced Tree (AVL) — height constraint:
  Balance factor = height(left) - height(right)
  Must be in {-1, 0, 1} for every node
  Rotation on violation:
    Left-Left: right rotate
    Right-Right: left rotate
    Left-Right: left rotate then right rotate
    Right-Left: right rotate then left rotate

Red-Black Tree (less strict, used in Java TreeMap):
  Properties:
  1. Every node is red or black
  2. Root is black
  3. Red nodes have black children
  4. All paths from node to null have same # black nodes
  → Height ≤ 2*log(n+1) guaranteed

Tree Traversals:
  Pre-order:  root → left → right   [1,2,4,5,3,6]
  In-order:   left → root → right   [4,2,5,1,3,6] ← sorted for BST!
  Post-order: left → right → root   [4,5,2,6,3,1]
  Level-order: BFS level by level   [1,2,3,4,5,6]
\end{verbatim}

\textbf{Complexity:}

\begin{longtable}[]{@{}
  >{\raggedright\arraybackslash}p{(\columnwidth - 6\tabcolsep) * \real{0.1514}}
  >{\raggedright\arraybackslash}p{(\columnwidth - 6\tabcolsep) * \real{0.2611}}
  >{\raggedright\arraybackslash}p{(\columnwidth - 6\tabcolsep) * \real{0.3046}}
  >{\raggedright\arraybackslash}p{(\columnwidth - 6\tabcolsep) * \real{0.2829}}@{}}
\toprule\noalign{}
\rowcolor{acc!16}
\begin{minipage}[b]{\linewidth}\raggedright
Op
\end{minipage} & \begin{minipage}[b]{\linewidth}\raggedright
Balanced BST
\end{minipage} & \begin{minipage}[b]{\linewidth}\raggedright
Unbalanced BST
\end{minipage} & \begin{minipage}[b]{\linewidth}\raggedright
AVL / RB Tree
\end{minipage} \\
\midrule\noalign{}
\endhead
\bottomrule\noalign{}
\endlastfoot
Search & O(log n) & O(n) & O(log n) \\
\rowcolor{zebra}
Insert & O(log n) & O(n) & O(log n) \\
Delete & O(log n) & O(n) & O(log n) \\
\rowcolor{zebra}
Min/Max & O(log n) & O(n) & O(log n) \\
In-order & O(n) & O(n) & O(n) \\
\end{longtable}

\textbf{When to use:}

\begin{itemize}
\tightlist
\item
  BST: sorted data with dynamic insertions/deletions
\item
  AVL/RB: guaranteed O(log n) when data can be adversarial
\item
  Any tree: hierarchical relationships, recursive decomposition
\end{itemize}

\textbf{Key Python patterns:}

\setcodelang{Python}

\begin{Shaded}
\begin{Highlighting}[]
\KeywordTok{class}\NormalTok{ TreeNode:}
    \KeywordTok{def} \FunctionTok{\_\_init\_\_}\NormalTok{(}\VariableTok{self}\NormalTok{, val}\OperatorTok{=}\DecValTok{0}\NormalTok{, left}\OperatorTok{=}\VariableTok{None}\NormalTok{, right}\OperatorTok{=}\VariableTok{None}\NormalTok{):}
        \VariableTok{self}\NormalTok{.val }\OperatorTok{=}\NormalTok{ val}
        \VariableTok{self}\NormalTok{.left }\OperatorTok{=}\NormalTok{ left}
        \VariableTok{self}\NormalTok{.right }\OperatorTok{=}\NormalTok{ right}

\CommentTok{\# In{-}order traversal (recursive)}
\KeywordTok{def}\NormalTok{ inorder(root):}
    \ControlFlowTok{if} \KeywordTok{not}\NormalTok{ root: }\ControlFlowTok{return}\NormalTok{ []}
    \ControlFlowTok{return}\NormalTok{ inorder(root.left) }\OperatorTok{+}\NormalTok{ [root.val] }\OperatorTok{+}\NormalTok{ inorder(root.right)}

\CommentTok{\# In{-}order traversal (iterative — important for interviews)}
\KeywordTok{def}\NormalTok{ inorder\_iter(root):}
\NormalTok{    result, stack }\OperatorTok{=}\NormalTok{ [], []}
\NormalTok{    curr }\OperatorTok{=}\NormalTok{ root}
    \ControlFlowTok{while}\NormalTok{ curr }\KeywordTok{or}\NormalTok{ stack:}
        \ControlFlowTok{while}\NormalTok{ curr:}
\NormalTok{            stack.append(curr)}
\NormalTok{            curr }\OperatorTok{=}\NormalTok{ curr.left}
\NormalTok{        curr }\OperatorTok{=}\NormalTok{ stack.pop()}
\NormalTok{        result.append(curr.val)}
\NormalTok{        curr }\OperatorTok{=}\NormalTok{ curr.right}
    \ControlFlowTok{return}\NormalTok{ result}

\CommentTok{\# Height of tree}
\KeywordTok{def}\NormalTok{ height(root):}
    \ControlFlowTok{if} \KeywordTok{not}\NormalTok{ root: }\ControlFlowTok{return} \DecValTok{0}
    \ControlFlowTok{return} \DecValTok{1} \OperatorTok{+} \BuiltInTok{max}\NormalTok{(height(root.left), height(root.right))}

\CommentTok{\# Check if BST is valid}
\KeywordTok{def}\NormalTok{ is\_valid\_bst(root, lo}\OperatorTok{=}\BuiltInTok{float}\NormalTok{(}\StringTok{\textquotesingle{}{-}inf\textquotesingle{}}\NormalTok{), hi}\OperatorTok{=}\BuiltInTok{float}\NormalTok{(}\StringTok{\textquotesingle{}inf\textquotesingle{}}\NormalTok{)):}
    \ControlFlowTok{if} \KeywordTok{not}\NormalTok{ root: }\ControlFlowTok{return} \VariableTok{True}
    \ControlFlowTok{if} \KeywordTok{not}\NormalTok{ (lo }\OperatorTok{\textless{}}\NormalTok{ root.val }\OperatorTok{\textless{}}\NormalTok{ hi): }\ControlFlowTok{return} \VariableTok{False}
    \ControlFlowTok{return}\NormalTok{ (is\_valid\_bst(root.left, lo, root.val) }\KeywordTok{and}
\NormalTok{            is\_valid\_bst(root.right, root.val, hi))}

\CommentTok{\# LCA (Lowest Common Ancestor) in BST}
\KeywordTok{def}\NormalTok{ lca\_bst(root, p, q):}
    \ControlFlowTok{if}\NormalTok{ p.val }\OperatorTok{\textless{}}\NormalTok{ root.val }\KeywordTok{and}\NormalTok{ q.val }\OperatorTok{\textless{}}\NormalTok{ root.val:}
        \ControlFlowTok{return}\NormalTok{ lca\_bst(root.left, p, q)}
    \ControlFlowTok{if}\NormalTok{ p.val }\OperatorTok{\textgreater{}}\NormalTok{ root.val }\KeywordTok{and}\NormalTok{ q.val }\OperatorTok{\textgreater{}}\NormalTok{ root.val:}
        \ControlFlowTok{return}\NormalTok{ lca\_bst(root.right, p, q)}
    \ControlFlowTok{return}\NormalTok{ root}

\CommentTok{\# Level order traversal}
\ImportTok{from}\NormalTok{ collections }\ImportTok{import}\NormalTok{ deque}
\KeywordTok{def}\NormalTok{ level\_order(root):}
    \ControlFlowTok{if} \KeywordTok{not}\NormalTok{ root: }\ControlFlowTok{return}\NormalTok{ []}
\NormalTok{    result, q }\OperatorTok{=}\NormalTok{ [], deque([root])}
    \ControlFlowTok{while}\NormalTok{ q:}
\NormalTok{        level }\OperatorTok{=}\NormalTok{ []}
        \ControlFlowTok{for}\NormalTok{ \_ }\KeywordTok{in} \BuiltInTok{range}\NormalTok{(}\BuiltInTok{len}\NormalTok{(q)):}
\NormalTok{            node }\OperatorTok{=}\NormalTok{ q.popleft()}
\NormalTok{            level.append(node.val)}
            \ControlFlowTok{if}\NormalTok{ node.left: q.append(node.left)}
            \ControlFlowTok{if}\NormalTok{ node.right: q.append(node.right)}
\NormalTok{        result.append(level)}
    \ControlFlowTok{return}\NormalTok{ result}
\end{Highlighting}
\end{Shaded}

\textbf{Interview takeaway:} \textbf{BST in-order = sorted. Use recursion with min/max bounds for BST validation, not just left \textless{} root \textless{} right (common bug). Diameter of tree = longest path through any node = max(left\_depth + right\_depth) over all nodes.}

\textbf{Real interview use-case:} Validate BST, lowest common ancestor, diameter of binary tree, serialize/deserialize tree, path sum, kth smallest in BST.

\subsection{Heaps / Priority Queues}\label{heaps--priority-queues}

\textbf{What it is:} A complete binary tree stored as an array, where each parent satisfies the heap property (min-heap: parent \(\leq\) children; max-heap: parent \(\geq\) children). Top element is always the min (or max).

\textbf{Internal mechanics:}

\setcodelang{}

\begin{verbatim}
Min-Heap (array representation):
         1
        / \
       3   5
      / \ / \
     7  4 8  9

Array: [1, 3, 5, 7, 4, 8, 9]
Index:  0  1  2  3  4  5  6

Index relationships:
  Parent(i)      = (i-1) // 2
  Left child(i)  = 2*i + 1
  Right child(i) = 2*i + 2

Sift Up (after insert at end):
  Insert 2: append → [1,3,5,7,4,8,9,2]
  Compare 2 with parent (index 3): arr[3]=7 > 2 → swap
  → [1,3,5,2,4,8,9,7]
  Compare 2 with parent (index 1): arr[1]=3 > 2 → swap
  → [1,2,5,3,4,8,9,7]
  Compare 2 with parent (index 0): arr[0]=1 < 2 → stop ✓

Sift Down (after extracting root — put last element at root):
  Extract 1: put 7 at root → [7,2,5,3,4,8,9]
  Compare 7 with children 2,5: min child=2 → swap
  → [2,7,5,3,4,8,9]
  Compare 7 with children 3,4: min child=3 → swap
  → [2,3,5,7,4,8,9]
  7 has no children → stop ✓

Build Heap (heapify) from arbitrary array — O(n):
  Start sifting down from last non-leaf (n//2 - 1) to root
  Counterintuitive: O(n) not O(n log n) because most nodes are near leaves
\end{verbatim}

\textbf{Complexity:}

\begin{longtable}[]{@{}
  >{\raggedright\arraybackslash}p{(\columnwidth - 2\tabcolsep) * \real{0.5238}}
  >{\raggedright\arraybackslash}p{(\columnwidth - 2\tabcolsep) * \real{0.4762}}@{}}
\toprule\noalign{}
\rowcolor{acc!16}
\begin{minipage}[b]{\linewidth}\raggedright
Op
\end{minipage} & \begin{minipage}[b]{\linewidth}\raggedright
Time
\end{minipage} \\
\midrule\noalign{}
\endhead
\bottomrule\noalign{}
\endlastfoot
Get min/max & O(1) \\
\rowcolor{zebra}
Insert & O(log n) \\
Extract min & O(log n) \\
\rowcolor{zebra}
Build heap & O(n) \\
Search & O(n) \\
\rowcolor{zebra}
Heapsort & O(n log n) \\
\end{longtable}

\textbf{When to use:}

\begin{itemize}
\tightlist
\item
  Get the min/max repeatedly (priority queues)
\item
  Top K elements (use heap of size K)
\item
  Merge K sorted lists
\item
  Dijkstra\textquotesingle s shortest path
\item
  Median maintenance (two heaps)
\end{itemize}

\textbf{Key Python patterns:}

\setcodelang{Python}

\begin{Shaded}
\begin{Highlighting}[]
\ImportTok{import}\NormalTok{ heapq}

\CommentTok{\# Min{-}heap (default)}
\NormalTok{heap }\OperatorTok{=}\NormalTok{ []}
\NormalTok{heapq.heappush(heap, }\DecValTok{3}\NormalTok{)}
\NormalTok{heapq.heappush(heap, }\DecValTok{1}\NormalTok{)}
\NormalTok{heapq.heappush(heap, }\DecValTok{2}\NormalTok{)}
\NormalTok{heapq.heappop(heap)    }\CommentTok{\# returns 1}

\CommentTok{\# Max{-}heap: negate values}
\NormalTok{max\_heap }\OperatorTok{=}\NormalTok{ []}
\NormalTok{heapq.heappush(max\_heap, }\OperatorTok{{-}}\DecValTok{3}\NormalTok{)}
\NormalTok{heapq.heappush(max\_heap, }\OperatorTok{{-}}\DecValTok{1}\NormalTok{)}
\OperatorTok{{-}}\NormalTok{heapq.heappop(max\_heap)   }\CommentTok{\# returns 3}

\CommentTok{\# Build heap in O(n)}
\NormalTok{nums }\OperatorTok{=}\NormalTok{ [}\DecValTok{3}\NormalTok{, }\DecValTok{1}\NormalTok{, }\DecValTok{4}\NormalTok{, }\DecValTok{1}\NormalTok{, }\DecValTok{5}\NormalTok{, }\DecValTok{9}\NormalTok{]}
\NormalTok{heapq.heapify(nums)}

\CommentTok{\# Top K elements}
\KeywordTok{def}\NormalTok{ top\_k(nums, k):}
    \ControlFlowTok{return}\NormalTok{ heapq.nlargest(k, nums)   }\CommentTok{\# O(n log k)}

\CommentTok{\# Top K using min{-}heap of size k (space efficient)}
\KeywordTok{def}\NormalTok{ top\_k\_heap(nums, k):}
\NormalTok{    heap }\OperatorTok{=}\NormalTok{ []}
    \ControlFlowTok{for}\NormalTok{ n }\KeywordTok{in}\NormalTok{ nums:}
\NormalTok{        heapq.heappush(heap, n)}
        \ControlFlowTok{if} \BuiltInTok{len}\NormalTok{(heap) }\OperatorTok{\textgreater{}}\NormalTok{ k:}
\NormalTok{            heapq.heappop(heap)}
    \ControlFlowTok{return} \BuiltInTok{sorted}\NormalTok{(heap, reverse}\OperatorTok{=}\VariableTok{True}\NormalTok{)}

\CommentTok{\# Merge K sorted lists}
\KeywordTok{def}\NormalTok{ merge\_k\_lists(lists):}
\NormalTok{    heap }\OperatorTok{=}\NormalTok{ []}
    \ControlFlowTok{for}\NormalTok{ i, lst }\KeywordTok{in} \BuiltInTok{enumerate}\NormalTok{(lists):}
        \ControlFlowTok{if}\NormalTok{ lst:}
\NormalTok{            heapq.heappush(heap, (lst[}\DecValTok{0}\NormalTok{], i, }\DecValTok{0}\NormalTok{))}
\NormalTok{    result }\OperatorTok{=}\NormalTok{ []}
    \ControlFlowTok{while}\NormalTok{ heap:}
\NormalTok{        val, i, j }\OperatorTok{=}\NormalTok{ heapq.heappop(heap)}
\NormalTok{        result.append(val)}
        \ControlFlowTok{if}\NormalTok{ j }\OperatorTok{+} \DecValTok{1} \OperatorTok{\textless{}} \BuiltInTok{len}\NormalTok{(lists[i]):}
\NormalTok{            heapq.heappush(heap, (lists[i][j}\OperatorTok{+}\DecValTok{1}\NormalTok{], i, j}\OperatorTok{+}\DecValTok{1}\NormalTok{))}
    \ControlFlowTok{return}\NormalTok{ result}

\CommentTok{\# Median maintenance (two heaps)}
\KeywordTok{class}\NormalTok{ MedianFinder:}
    \KeywordTok{def} \FunctionTok{\_\_init\_\_}\NormalTok{(}\VariableTok{self}\NormalTok{):}
        \VariableTok{self}\NormalTok{.lo }\OperatorTok{=}\NormalTok{ []   }\CommentTok{\# max{-}heap (negate)}
        \VariableTok{self}\NormalTok{.hi }\OperatorTok{=}\NormalTok{ []   }\CommentTok{\# min{-}heap}
    \KeywordTok{def}\NormalTok{ add\_num(}\VariableTok{self}\NormalTok{, num):}
\NormalTok{        heapq.heappush(}\VariableTok{self}\NormalTok{.lo, }\OperatorTok{{-}}\NormalTok{num)}
\NormalTok{        heapq.heappush(}\VariableTok{self}\NormalTok{.hi, }\OperatorTok{{-}}\NormalTok{heapq.heappop(}\VariableTok{self}\NormalTok{.lo))}
        \ControlFlowTok{if} \BuiltInTok{len}\NormalTok{(}\VariableTok{self}\NormalTok{.hi) }\OperatorTok{\textgreater{}} \BuiltInTok{len}\NormalTok{(}\VariableTok{self}\NormalTok{.lo):}
\NormalTok{            heapq.heappush(}\VariableTok{self}\NormalTok{.lo, }\OperatorTok{{-}}\NormalTok{heapq.heappop(}\VariableTok{self}\NormalTok{.hi))}
    \KeywordTok{def}\NormalTok{ find\_median(}\VariableTok{self}\NormalTok{):}
        \ControlFlowTok{if} \BuiltInTok{len}\NormalTok{(}\VariableTok{self}\NormalTok{.lo) }\OperatorTok{\textgreater{}} \BuiltInTok{len}\NormalTok{(}\VariableTok{self}\NormalTok{.hi): }\ControlFlowTok{return} \OperatorTok{{-}}\VariableTok{self}\NormalTok{.lo[}\DecValTok{0}\NormalTok{]}
        \ControlFlowTok{return}\NormalTok{ (}\OperatorTok{{-}}\VariableTok{self}\NormalTok{.lo[}\DecValTok{0}\NormalTok{] }\OperatorTok{+} \VariableTok{self}\NormalTok{.hi[}\DecValTok{0}\NormalTok{]) }\OperatorTok{/} \FloatTok{2.0}
\end{Highlighting}
\end{Shaded}

\textbf{Interview takeaway:} \textbf{Python only has min-heap via heapq; negate values for max-heap. \texttt{heapify} is O(n). Top K smallest: use max-heap of size K and push out the max when size exceeds K. Top K largest: use min-heap of size K.}

\textbf{Real interview use-case:} Kth largest element, merge K sorted lists, task scheduler, find median from data stream, Dijkstra\textquotesingle s.

\subsection{Tries (Prefix Trees)}\label{tries-prefix-trees}

\textbf{What it is:} A tree where each path from root to leaf spells a word, and each node represents one character. Shared prefixes share nodes.

\textbf{Internal mechanics:}

\setcodelang{}

\begin{verbatim}
Trie storing: ["cat", "car", "card", "care", "bat"]

          root
         /    \
        c      b
        |      |
        a      a
       / \      \
      t   r      t
     (*)  |      (*)
         / \
        d   e
       (*)  (*)

(*) = is_end marker

Node structure:
  children: dict[char → TrieNode]  (or array[26])
  is_end: bool

Insert "card":
  root → c (create if absent)
       → a (create if absent)
       → r (exists, reuse)
       → d (create)
       mark d.is_end = True

Search "care":
  root → c → a → r → e → is_end? True ✓

StartsWith "ca":
  root → c → a → exists? True ✓ (prefix search)

Space: each node uses O(ALPHA) where ALPHA=26 for lowercase
  n strings of avg length m: O(n*m) nodes worst case
  (better than storing all strings if many shared prefixes)
\end{verbatim}

\textbf{Complexity (m = length of key):}

\begin{longtable}[]{@{}
  >{\raggedright\arraybackslash}p{(\columnwidth - 4\tabcolsep) * \real{0.4914}}
  >{\raggedright\arraybackslash}p{(\columnwidth - 4\tabcolsep) * \real{0.2260}}
  >{\raggedright\arraybackslash}p{(\columnwidth - 4\tabcolsep) * \real{0.2826}}@{}}
\toprule\noalign{}
\rowcolor{acc!16}
\begin{minipage}[b]{\linewidth}\raggedright
Op
\end{minipage} & \begin{minipage}[b]{\linewidth}\raggedright
Time
\end{minipage} & \begin{minipage}[b]{\linewidth}\raggedright
Space
\end{minipage} \\
\midrule\noalign{}
\endhead
\bottomrule\noalign{}
\endlastfoot
Insert & O(m) & O(m) \\
\rowcolor{zebra}
Search & O(m) & O(1) \\
StartsWith & O(m) & O(1) \\
\rowcolor{zebra}
Delete & O(m) & O(1) \\
\end{longtable}

\textbf{When to use:}

\begin{itemize}
\tightlist
\item
  Autocomplete / prefix search
\item
  Spell checking
\item
  IP routing tables
\item
  Word search in a grid
\end{itemize}

\textbf{When NOT to use:}

\begin{itemize}
\tightlist
\item
  Simple key-value lookup (use hash map)
\item
  Keys are not strings or have no shared prefixes
\end{itemize}

\textbf{Key Python patterns:}

\setcodelang{Python}

\begin{Shaded}
\begin{Highlighting}[]
\KeywordTok{class}\NormalTok{ TrieNode:}
    \KeywordTok{def} \FunctionTok{\_\_init\_\_}\NormalTok{(}\VariableTok{self}\NormalTok{):}
        \VariableTok{self}\NormalTok{.children }\OperatorTok{=}\NormalTok{ \{\}}
        \VariableTok{self}\NormalTok{.is\_end }\OperatorTok{=} \VariableTok{False}

\KeywordTok{class}\NormalTok{ Trie:}
    \KeywordTok{def} \FunctionTok{\_\_init\_\_}\NormalTok{(}\VariableTok{self}\NormalTok{):}
        \VariableTok{self}\NormalTok{.root }\OperatorTok{=}\NormalTok{ TrieNode()}

    \KeywordTok{def}\NormalTok{ insert(}\VariableTok{self}\NormalTok{, word: }\BuiltInTok{str}\NormalTok{) }\OperatorTok{{-}\textgreater{}} \VariableTok{None}\NormalTok{:}
\NormalTok{        node }\OperatorTok{=} \VariableTok{self}\NormalTok{.root}
        \ControlFlowTok{for}\NormalTok{ ch }\KeywordTok{in}\NormalTok{ word:}
            \ControlFlowTok{if}\NormalTok{ ch }\KeywordTok{not} \KeywordTok{in}\NormalTok{ node.children:}
\NormalTok{                node.children[ch] }\OperatorTok{=}\NormalTok{ TrieNode()}
\NormalTok{            node }\OperatorTok{=}\NormalTok{ node.children[ch]}
\NormalTok{        node.is\_end }\OperatorTok{=} \VariableTok{True}

    \KeywordTok{def}\NormalTok{ search(}\VariableTok{self}\NormalTok{, word: }\BuiltInTok{str}\NormalTok{) }\OperatorTok{{-}\textgreater{}} \BuiltInTok{bool}\NormalTok{:}
\NormalTok{        node }\OperatorTok{=} \VariableTok{self}\NormalTok{.root}
        \ControlFlowTok{for}\NormalTok{ ch }\KeywordTok{in}\NormalTok{ word:}
            \ControlFlowTok{if}\NormalTok{ ch }\KeywordTok{not} \KeywordTok{in}\NormalTok{ node.children:}
                \ControlFlowTok{return} \VariableTok{False}
\NormalTok{            node }\OperatorTok{=}\NormalTok{ node.children[ch]}
        \ControlFlowTok{return}\NormalTok{ node.is\_end}

    \KeywordTok{def}\NormalTok{ starts\_with(}\VariableTok{self}\NormalTok{, prefix: }\BuiltInTok{str}\NormalTok{) }\OperatorTok{{-}\textgreater{}} \BuiltInTok{bool}\NormalTok{:}
\NormalTok{        node }\OperatorTok{=} \VariableTok{self}\NormalTok{.root}
        \ControlFlowTok{for}\NormalTok{ ch }\KeywordTok{in}\NormalTok{ prefix:}
            \ControlFlowTok{if}\NormalTok{ ch }\KeywordTok{not} \KeywordTok{in}\NormalTok{ node.children:}
                \ControlFlowTok{return} \VariableTok{False}
\NormalTok{            node }\OperatorTok{=}\NormalTok{ node.children[ch]}
        \ControlFlowTok{return} \VariableTok{True}

\CommentTok{\# Word search II (trie + DFS backtracking)}
\CommentTok{\# Build trie from word list, then DFS on board}
\CommentTok{\# When we reach a trie end node, we found a word}
\end{Highlighting}
\end{Shaded}

\textbf{Interview takeaway:} \textbf{Trie is the answer to "find all words with prefix X" or "word search on a grid with a word list". For word search II (board + word list), build a trie from words and run DFS on the board --- this prunes early when no prefix matches.}

\textbf{Real interview use-case:} Implement trie, word search II, replace words, design search autocomplete, map sum pairs.

\subsection{Graphs}\label{graphs}

\textbf{What it is:} A set of vertices (nodes) connected by edges. Can be directed/undirected, weighted/unweighted, cyclic/acyclic.

\textbf{Internal mechanics:}

\setcodelang{}

\begin{verbatim}
Graph: 4 vertices, 4 edges (undirected)
  0 --- 1
  |   / |
  |  /  |
  | /   |
  2 --- 3

Adjacency List (preferred for sparse graphs):
  {
    0: [1, 2],
    1: [0, 2, 3],
    2: [0, 1, 3],
    3: [1, 2]
  }
  Space: O(V + E)
  Check edge (u,v): O(degree(u))
  Get neighbors: O(degree(u))

Adjacency Matrix (preferred for dense graphs):
      0  1  2  3
  0 [ 0  1  1  0 ]
  1 [ 1  0  1  1 ]
  2 [ 1  1  0  1 ]
  3 [ 0  1  1  0 ]
  Space: O(V²)
  Check edge (u,v): O(1)
  Get neighbors: O(V)

Weighted adjacency list:
  {0: [(1, 4), (2, 1)], ...}  ← (neighbor, weight)

Types:
  Undirected: edges go both ways
  Directed (Digraph): edges have direction
  DAG: Directed Acyclic Graph (used in topological sort)
  Tree: connected undirected graph with V-1 edges
\end{verbatim}

\textbf{Complexity:}

\begin{longtable}[]{@{}
  >{\raggedright\arraybackslash}p{(\columnwidth - 4\tabcolsep) * \real{0.4198}}
  >{\raggedright\arraybackslash}p{(\columnwidth - 4\tabcolsep) * \real{0.2963}}
  >{\raggedright\arraybackslash}p{(\columnwidth - 4\tabcolsep) * \real{0.2840}}@{}}
\toprule\noalign{}
\rowcolor{acc!16}
\begin{minipage}[b]{\linewidth}\raggedright
Op
\end{minipage} & \begin{minipage}[b]{\linewidth}\raggedright
Adj List
\end{minipage} & \begin{minipage}[b]{\linewidth}\raggedright
Adj Matrix
\end{minipage} \\
\midrule\noalign{}
\endhead
\bottomrule\noalign{}
\endlastfoot
Space & O(V+E) & O(V\(^2\)) \\
\rowcolor{zebra}
Add vertex & O(1) & O(V\(^2\)) \\
Add edge & O(1) & O(1) \\
\rowcolor{zebra}
Remove edge & O(E) & O(1) \\
Check edge (u,v) & O(degree) & O(1) \\
\rowcolor{zebra}
Get all neighbors & O(degree) & O(V) \\
BFS/DFS & O(V+E) & O(V\(^2\)) \\
\rowcolor{zebra}
Dijkstra & O((V+E)logV) & O(V\(^2\)logV) \\
\end{longtable}

\textbf{When to use:}

\begin{itemize}
\tightlist
\item
  Adj List: default, sparse graphs, BFS/DFS heavy
\item
  Adj Matrix: dense graphs, need O(1) edge existence check, Floyd-Warshall
\end{itemize}

\textbf{Key Python patterns:}

\setcodelang{Python}

\begin{Shaded}
\begin{Highlighting}[]
\ImportTok{from}\NormalTok{ collections }\ImportTok{import}\NormalTok{ defaultdict, deque}

\CommentTok{\# Build graph}
\NormalTok{graph }\OperatorTok{=}\NormalTok{ defaultdict(}\BuiltInTok{list}\NormalTok{)}
\ControlFlowTok{for}\NormalTok{ u, v }\KeywordTok{in}\NormalTok{ edges:}
\NormalTok{    graph[u].append(v)}
\NormalTok{    graph[v].append(u)  }\CommentTok{\# undirected}

\CommentTok{\# BFS (shortest path in unweighted graph)}
\KeywordTok{def}\NormalTok{ bfs(graph, start, end):}
\NormalTok{    visited }\OperatorTok{=}\NormalTok{ \{start\}}
\NormalTok{    q }\OperatorTok{=}\NormalTok{ deque([(start, }\DecValTok{0}\NormalTok{)])}
    \ControlFlowTok{while}\NormalTok{ q:}
\NormalTok{        node, dist }\OperatorTok{=}\NormalTok{ q.popleft()}
        \ControlFlowTok{if}\NormalTok{ node }\OperatorTok{==}\NormalTok{ end: }\ControlFlowTok{return}\NormalTok{ dist}
        \ControlFlowTok{for}\NormalTok{ nei }\KeywordTok{in}\NormalTok{ graph[node]:}
            \ControlFlowTok{if}\NormalTok{ nei }\KeywordTok{not} \KeywordTok{in}\NormalTok{ visited:}
\NormalTok{                visited.add(nei)}
\NormalTok{                q.append((nei, dist }\OperatorTok{+} \DecValTok{1}\NormalTok{))}
    \ControlFlowTok{return} \OperatorTok{{-}}\DecValTok{1}

\CommentTok{\# DFS (iterative)}
\KeywordTok{def}\NormalTok{ dfs(graph, start):}
\NormalTok{    visited }\OperatorTok{=} \BuiltInTok{set}\NormalTok{()}
\NormalTok{    stack }\OperatorTok{=}\NormalTok{ [start]}
    \ControlFlowTok{while}\NormalTok{ stack:}
\NormalTok{        node }\OperatorTok{=}\NormalTok{ stack.pop()}
        \ControlFlowTok{if}\NormalTok{ node }\KeywordTok{in}\NormalTok{ visited: }\ControlFlowTok{continue}
\NormalTok{        visited.add(node)}
        \ControlFlowTok{for}\NormalTok{ nei }\KeywordTok{in}\NormalTok{ graph[node]:}
\NormalTok{            stack.append(nei)}

\CommentTok{\# DFS (recursive)}
\KeywordTok{def}\NormalTok{ dfs\_rec(node, visited, graph):}
\NormalTok{    visited.add(node)}
    \ControlFlowTok{for}\NormalTok{ nei }\KeywordTok{in}\NormalTok{ graph[node]:}
        \ControlFlowTok{if}\NormalTok{ nei }\KeywordTok{not} \KeywordTok{in}\NormalTok{ visited:}
\NormalTok{            dfs\_rec(nei, visited, graph)}

\CommentTok{\# Topological sort (Kahn\textquotesingle{}s algorithm, BFS{-}based)}
\KeywordTok{def}\NormalTok{ topo\_sort(n, edges):}
\NormalTok{    graph }\OperatorTok{=}\NormalTok{ defaultdict(}\BuiltInTok{list}\NormalTok{)}
\NormalTok{    in\_degree }\OperatorTok{=}\NormalTok{ [}\DecValTok{0}\NormalTok{] }\OperatorTok{*}\NormalTok{ n}
    \ControlFlowTok{for}\NormalTok{ u, v }\KeywordTok{in}\NormalTok{ edges:}
\NormalTok{        graph[u].append(v)}
\NormalTok{        in\_degree[v] }\OperatorTok{+=} \DecValTok{1}
\NormalTok{    q }\OperatorTok{=}\NormalTok{ deque(i }\ControlFlowTok{for}\NormalTok{ i }\KeywordTok{in} \BuiltInTok{range}\NormalTok{(n) }\ControlFlowTok{if}\NormalTok{ in\_degree[i] }\OperatorTok{==} \DecValTok{0}\NormalTok{)}
\NormalTok{    order }\OperatorTok{=}\NormalTok{ []}
    \ControlFlowTok{while}\NormalTok{ q:}
\NormalTok{        node }\OperatorTok{=}\NormalTok{ q.popleft()}
\NormalTok{        order.append(node)}
        \ControlFlowTok{for}\NormalTok{ nei }\KeywordTok{in}\NormalTok{ graph[node]:}
\NormalTok{            in\_degree[nei] }\OperatorTok{{-}=} \DecValTok{1}
            \ControlFlowTok{if}\NormalTok{ in\_degree[nei] }\OperatorTok{==} \DecValTok{0}\NormalTok{:}
\NormalTok{                q.append(nei)}
    \ControlFlowTok{return}\NormalTok{ order }\ControlFlowTok{if} \BuiltInTok{len}\NormalTok{(order) }\OperatorTok{==}\NormalTok{ n }\ControlFlowTok{else}\NormalTok{ []  }\CommentTok{\# empty = cycle exists}

\CommentTok{\# Dijkstra\textquotesingle{}s (single source shortest path)}
\ImportTok{import}\NormalTok{ heapq}
\KeywordTok{def}\NormalTok{ dijkstra(graph, start):}
\NormalTok{    dist }\OperatorTok{=}\NormalTok{ \{node: }\BuiltInTok{float}\NormalTok{(}\StringTok{\textquotesingle{}inf\textquotesingle{}}\NormalTok{) }\ControlFlowTok{for}\NormalTok{ node }\KeywordTok{in}\NormalTok{ graph\}}
\NormalTok{    dist[start] }\OperatorTok{=} \DecValTok{0}
\NormalTok{    heap }\OperatorTok{=}\NormalTok{ [(}\DecValTok{0}\NormalTok{, start)]}
    \ControlFlowTok{while}\NormalTok{ heap:}
\NormalTok{        d, node }\OperatorTok{=}\NormalTok{ heapq.heappop(heap)}
        \ControlFlowTok{if}\NormalTok{ d }\OperatorTok{\textgreater{}}\NormalTok{ dist[node]: }\ControlFlowTok{continue}
        \ControlFlowTok{for}\NormalTok{ nei, weight }\KeywordTok{in}\NormalTok{ graph[node]:}
\NormalTok{            new\_d }\OperatorTok{=}\NormalTok{ d }\OperatorTok{+}\NormalTok{ weight}
            \ControlFlowTok{if}\NormalTok{ new\_d }\OperatorTok{\textless{}}\NormalTok{ dist[nei]:}
\NormalTok{                dist[nei] }\OperatorTok{=}\NormalTok{ new\_d}
\NormalTok{                heapq.heappush(heap, (new\_d, nei))}
    \ControlFlowTok{return}\NormalTok{ dist}
\end{Highlighting}
\end{Shaded}

\textbf{Interview takeaway:} \textbf{BFS for shortest path in unweighted graph. Dijkstra for weighted (non-negative). Detect cycle in directed graph: DFS with in-stack tracking. Topological sort only valid for DAGs. "Number of islands" = number of DFS/BFS calls on unvisited cells.}

\textbf{Real interview use-case:} Number of islands, clone graph, course schedule, network delay time, Pacific Atlantic water flow, word ladder.

\subsection{Union-Find / Disjoint Set}\label{union-find--disjoint-set}

\textbf{What it is:} Data structure to track which elements belong to the same connected component. Supports two operations: \texttt{find} (which component?) and \texttt{union} (merge components).

\textbf{Internal mechanics:}

\setcodelang{}

\begin{verbatim}
Initial state (6 elements, each its own component):
  parent = [0, 1, 2, 3, 4, 5]
  rank   = [0, 0, 0, 0, 0, 0]

Union(0, 1): find roots → 0,1; rank equal → make 1 child of 0
  parent = [0, 0, 2, 3, 4, 5]
  rank   = [1, 0, 0, 0, 0, 0]

Union(2, 3):
  parent = [0, 0, 2, 2, 4, 5]

Union(0, 2): find roots → 0, 2; rank[0]=1 > rank[2]=0 → 2 under 0
  parent = [0, 0, 0, 2, 4, 5]   ← 2's parent is now 0

Path Compression during find(3):
  find(3): parent[3]=2, parent[2]=0, parent[0]=0 → root=0
  Compress: set parent[3]=0, parent[2]=0 directly
  parent = [0, 0, 0, 0, 4, 5]   ← next find(3) is O(1)

After path compression and union by rank:
  find and union are effectively O(α(n)) ≈ O(1)
  where α = inverse Ackermann function
  α(n) < 5 for any practical n

Visual after several unions:
  Component 1: {0,1,2,3}  ← tree rooted at 0
  Component 2: {4}
  Component 3: {5}
       0
      /|\
     1 2 3   ← after path compression, flat tree
\end{verbatim}

\textbf{Complexity:}

\begin{longtable}[]{@{}
  >{\raggedright\arraybackslash}p{(\columnwidth - 6\tabcolsep) * \real{0.0986}}
  >{\raggedright\arraybackslash}p{(\columnwidth - 6\tabcolsep) * \real{0.1134}}
  >{\raggedright\arraybackslash}p{(\columnwidth - 6\tabcolsep) * \real{0.2949}}
  >{\raggedright\arraybackslash}p{(\columnwidth - 6\tabcolsep) * \real{0.4931}}@{}}
\toprule\noalign{}
\rowcolor{acc!16}
\begin{minipage}[b]{\linewidth}\raggedright
Op
\end{minipage} & \begin{minipage}[b]{\linewidth}\raggedright
Naive
\end{minipage} & \begin{minipage}[b]{\linewidth}\raggedright
Union by Rank
\end{minipage} & \begin{minipage}[b]{\linewidth}\raggedright
+ Path Compression
\end{minipage} \\
\midrule\noalign{}
\endhead
\bottomrule\noalign{}
\endlastfoot
Find & O(n) & O(log n) & O(\(\alpha\)(n)) \(\approx\) O(1) \\
\rowcolor{zebra}
Union & O(n) & O(log n) & O(\(\alpha\)(n)) \(\approx\) O(1) \\
Space & O(n) & O(n) & O(n) \\
\end{longtable}

\textbf{When to use:}

\begin{itemize}
\tightlist
\item
  Connected components (dynamic)
\item
  Detect cycles in undirected graph
\item
  Kruskal\textquotesingle s minimum spanning tree
\item
  Grouping problems ("same group as")
\end{itemize}

\textbf{Key Python patterns:}

\setcodelang{Python}

\begin{Shaded}
\begin{Highlighting}[]
\KeywordTok{class}\NormalTok{ UnionFind:}
    \KeywordTok{def} \FunctionTok{\_\_init\_\_}\NormalTok{(}\VariableTok{self}\NormalTok{, n):}
        \VariableTok{self}\NormalTok{.parent }\OperatorTok{=} \BuiltInTok{list}\NormalTok{(}\BuiltInTok{range}\NormalTok{(n))}
        \VariableTok{self}\NormalTok{.rank }\OperatorTok{=}\NormalTok{ [}\DecValTok{0}\NormalTok{] }\OperatorTok{*}\NormalTok{ n}
        \VariableTok{self}\NormalTok{.components }\OperatorTok{=}\NormalTok{ n}

    \KeywordTok{def}\NormalTok{ find(}\VariableTok{self}\NormalTok{, x):}
        \ControlFlowTok{if} \VariableTok{self}\NormalTok{.parent[x] }\OperatorTok{!=}\NormalTok{ x:}
            \VariableTok{self}\NormalTok{.parent[x] }\OperatorTok{=} \VariableTok{self}\NormalTok{.find(}\VariableTok{self}\NormalTok{.parent[x])  }\CommentTok{\# path compression}
        \ControlFlowTok{return} \VariableTok{self}\NormalTok{.parent[x]}

    \KeywordTok{def}\NormalTok{ union(}\VariableTok{self}\NormalTok{, x, y):}
\NormalTok{        rx, ry }\OperatorTok{=} \VariableTok{self}\NormalTok{.find(x), }\VariableTok{self}\NormalTok{.find(y)}
        \ControlFlowTok{if}\NormalTok{ rx }\OperatorTok{==}\NormalTok{ ry: }\ControlFlowTok{return} \VariableTok{False}  \CommentTok{\# already connected}
        \CommentTok{\# Union by rank}
        \ControlFlowTok{if} \VariableTok{self}\NormalTok{.rank[rx] }\OperatorTok{\textless{}} \VariableTok{self}\NormalTok{.rank[ry]:}
\NormalTok{            rx, ry }\OperatorTok{=}\NormalTok{ ry, rx}
        \VariableTok{self}\NormalTok{.parent[ry] }\OperatorTok{=}\NormalTok{ rx}
        \ControlFlowTok{if} \VariableTok{self}\NormalTok{.rank[rx] }\OperatorTok{==} \VariableTok{self}\NormalTok{.rank[ry]:}
            \VariableTok{self}\NormalTok{.rank[rx] }\OperatorTok{+=} \DecValTok{1}
        \VariableTok{self}\NormalTok{.components }\OperatorTok{{-}=} \DecValTok{1}
        \ControlFlowTok{return} \VariableTok{True}

    \KeywordTok{def}\NormalTok{ connected(}\VariableTok{self}\NormalTok{, x, y):}
        \ControlFlowTok{return} \VariableTok{self}\NormalTok{.find(x) }\OperatorTok{==} \VariableTok{self}\NormalTok{.find(y)}

\CommentTok{\# Example: count components in graph}
\KeywordTok{def}\NormalTok{ count\_components(n, edges):}
\NormalTok{    uf }\OperatorTok{=}\NormalTok{ UnionFind(n)}
    \ControlFlowTok{for}\NormalTok{ u, v }\KeywordTok{in}\NormalTok{ edges:}
\NormalTok{        uf.union(u, v)}
    \ControlFlowTok{return}\NormalTok{ uf.components}

\CommentTok{\# Detect cycle in undirected graph}
\KeywordTok{def}\NormalTok{ has\_cycle(n, edges):}
\NormalTok{    uf }\OperatorTok{=}\NormalTok{ UnionFind(n)}
    \ControlFlowTok{for}\NormalTok{ u, v }\KeywordTok{in}\NormalTok{ edges:}
        \ControlFlowTok{if} \KeywordTok{not}\NormalTok{ uf.union(u, v):  }\CommentTok{\# already in same component}
            \ControlFlowTok{return} \VariableTok{True}
    \ControlFlowTok{return} \VariableTok{False}

\CommentTok{\# Kruskal\textquotesingle{}s MST}
\KeywordTok{def}\NormalTok{ kruskal(n, edges):}
\NormalTok{    uf }\OperatorTok{=}\NormalTok{ UnionFind(n)}
\NormalTok{    edges.sort(key}\OperatorTok{=}\KeywordTok{lambda}\NormalTok{ e: e[}\DecValTok{2}\NormalTok{])  }\CommentTok{\# sort by weight}
\NormalTok{    mst\_cost }\OperatorTok{=} \DecValTok{0}
    \ControlFlowTok{for}\NormalTok{ u, v, w }\KeywordTok{in}\NormalTok{ edges:}
        \ControlFlowTok{if}\NormalTok{ uf.union(u, v):}
\NormalTok{            mst\_cost }\OperatorTok{+=}\NormalTok{ w}
    \ControlFlowTok{return}\NormalTok{ mst\_cost}
\end{Highlighting}
\end{Shaded}

\textbf{Interview takeaway:} \textbf{Always implement both path compression AND union by rank. Without both, you don\textquotesingle t get the near-O(1) guarantee. Union-Find cannot efficiently handle edge deletion --- if you need that, use Link-Cut Trees (not interview scope).}

\textbf{Real interview use-case:} Number of provinces, redundant connection, accounts merge, satisfiability of equality equations.

\subsection{LRU / LFU Cache Structures}\label{lru--lfu-cache-structures}

\textbf{What it is:}

\begin{itemize}
\tightlist
\item
  \textbf{LRU (Least Recently Used)}: Evicts the item not accessed for the longest time. O(1) get and put.
\item
  \textbf{LFU (Least Frequently Used)}: Evicts the item used fewest times (break ties by LRU). O(1) all operations.
\end{itemize}

\textbf{Internal mechanics:}

\setcodelang{}

\begin{verbatim}
LRU Cache (capacity=3):
  Internals: HashMap + Doubly Linked List
  Map: key → node (O(1) access)
  DLL: most recent at front, least recent at tail

  get(1): move node 1 to front
  put(4): if full, remove tail (LRU), insert at front, update map

  State after put(1), put(2), put(3):
    [3] ↔ [2] ↔ [1] ← tail (LRU)
     ↑
    head (MRU)

  get(1): move 1 to head
    [1] ↔ [3] ↔ [2] ← tail (LRU)

  put(4): evict tail (2), add 4 at head
    [4] ↔ [1] ↔ [3]

LFU Cache (capacity=3):
  Internals: HashMap(key→(val,freq)) + HashMap(freq→DLL of keys)
             + min_freq tracker
  On access: increment freq, move to freq+1 list
  On evict: remove LRU from min_freq list
  Trickier: when min_freq list empties, min_freq++
\end{verbatim}

\textbf{Complexity:}

\begin{longtable}[]{@{}
  >{\raggedright\arraybackslash}p{(\columnwidth - 4\tabcolsep) * \real{0.3333}}
  >{\raggedright\arraybackslash}p{(\columnwidth - 4\tabcolsep) * \real{0.3333}}
  >{\raggedright\arraybackslash}p{(\columnwidth - 4\tabcolsep) * \real{0.3333}}@{}}
\toprule\noalign{}
\rowcolor{acc!16}
\begin{minipage}[b]{\linewidth}\raggedright
Op
\end{minipage} & \begin{minipage}[b]{\linewidth}\raggedright
LRU
\end{minipage} & \begin{minipage}[b]{\linewidth}\raggedright
LFU
\end{minipage} \\
\midrule\noalign{}
\endhead
\bottomrule\noalign{}
\endlastfoot
get & O(1) & O(1) \\
\rowcolor{zebra}
put & O(1) & O(1) \\
\end{longtable}

\textbf{Key Python patterns:}

\setcodelang{Python}

\begin{Shaded}
\begin{Highlighting}[]
\ImportTok{from}\NormalTok{ collections }\ImportTok{import}\NormalTok{ OrderedDict}

\CommentTok{\# LRU using OrderedDict (move\_to\_end trick)}
\KeywordTok{class}\NormalTok{ LRUCache:}
    \KeywordTok{def} \FunctionTok{\_\_init\_\_}\NormalTok{(}\VariableTok{self}\NormalTok{, capacity: }\BuiltInTok{int}\NormalTok{):}
        \VariableTok{self}\NormalTok{.cap }\OperatorTok{=}\NormalTok{ capacity}
        \VariableTok{self}\NormalTok{.cache }\OperatorTok{=}\NormalTok{ OrderedDict()  }\CommentTok{\# key → value, ordered by access time}

    \KeywordTok{def}\NormalTok{ get(}\VariableTok{self}\NormalTok{, key: }\BuiltInTok{int}\NormalTok{) }\OperatorTok{{-}\textgreater{}} \BuiltInTok{int}\NormalTok{:}
        \ControlFlowTok{if}\NormalTok{ key }\KeywordTok{not} \KeywordTok{in} \VariableTok{self}\NormalTok{.cache: }\ControlFlowTok{return} \OperatorTok{{-}}\DecValTok{1}
        \VariableTok{self}\NormalTok{.cache.move\_to\_end(key)  }\CommentTok{\# mark as most recently used}
        \ControlFlowTok{return} \VariableTok{self}\NormalTok{.cache[key]}

    \KeywordTok{def}\NormalTok{ put(}\VariableTok{self}\NormalTok{, key: }\BuiltInTok{int}\NormalTok{, value: }\BuiltInTok{int}\NormalTok{) }\OperatorTok{{-}\textgreater{}} \VariableTok{None}\NormalTok{:}
        \ControlFlowTok{if}\NormalTok{ key }\KeywordTok{in} \VariableTok{self}\NormalTok{.cache:}
            \VariableTok{self}\NormalTok{.cache.move\_to\_end(key)}
        \VariableTok{self}\NormalTok{.cache[key] }\OperatorTok{=}\NormalTok{ value}
        \ControlFlowTok{if} \BuiltInTok{len}\NormalTok{(}\VariableTok{self}\NormalTok{.cache) }\OperatorTok{\textgreater{}} \VariableTok{self}\NormalTok{.cap:}
            \VariableTok{self}\NormalTok{.cache.popitem(last}\OperatorTok{=}\VariableTok{False}\NormalTok{)  }\CommentTok{\# remove LRU (first item)}

\CommentTok{\# LRU from scratch (HashMap + DLL) — what interviewers actually want}
\KeywordTok{class}\NormalTok{ DLLNode:}
    \KeywordTok{def} \FunctionTok{\_\_init\_\_}\NormalTok{(}\VariableTok{self}\NormalTok{, key}\OperatorTok{=}\DecValTok{0}\NormalTok{, val}\OperatorTok{=}\DecValTok{0}\NormalTok{):}
        \VariableTok{self}\NormalTok{.key }\OperatorTok{=}\NormalTok{ key}
        \VariableTok{self}\NormalTok{.val }\OperatorTok{=}\NormalTok{ val}
        \VariableTok{self}\NormalTok{.prev }\OperatorTok{=} \VariableTok{self}\NormalTok{.}\BuiltInTok{next} \OperatorTok{=} \VariableTok{None}

\KeywordTok{class}\NormalTok{ LRUCacheManual:}
    \KeywordTok{def} \FunctionTok{\_\_init\_\_}\NormalTok{(}\VariableTok{self}\NormalTok{, capacity):}
        \VariableTok{self}\NormalTok{.cap }\OperatorTok{=}\NormalTok{ capacity}
        \VariableTok{self}\NormalTok{.cache }\OperatorTok{=}\NormalTok{ \{\}}
        \VariableTok{self}\NormalTok{.head }\OperatorTok{=}\NormalTok{ DLLNode()   }\CommentTok{\# dummy head (MRU side)}
        \VariableTok{self}\NormalTok{.tail }\OperatorTok{=}\NormalTok{ DLLNode()   }\CommentTok{\# dummy tail (LRU side)}
        \VariableTok{self}\NormalTok{.head.}\BuiltInTok{next} \OperatorTok{=} \VariableTok{self}\NormalTok{.tail}
        \VariableTok{self}\NormalTok{.tail.prev }\OperatorTok{=} \VariableTok{self}\NormalTok{.head}

    \KeywordTok{def}\NormalTok{ \_remove(}\VariableTok{self}\NormalTok{, node):}
\NormalTok{        node.prev.}\BuiltInTok{next} \OperatorTok{=}\NormalTok{ node.}\BuiltInTok{next}
\NormalTok{        node.}\BuiltInTok{next}\NormalTok{.prev }\OperatorTok{=}\NormalTok{ node.prev}

    \KeywordTok{def}\NormalTok{ \_insert\_front(}\VariableTok{self}\NormalTok{, node):}
\NormalTok{        node.}\BuiltInTok{next} \OperatorTok{=} \VariableTok{self}\NormalTok{.head.}\BuiltInTok{next}
\NormalTok{        node.prev }\OperatorTok{=} \VariableTok{self}\NormalTok{.head}
        \VariableTok{self}\NormalTok{.head.}\BuiltInTok{next}\NormalTok{.prev }\OperatorTok{=}\NormalTok{ node}
        \VariableTok{self}\NormalTok{.head.}\BuiltInTok{next} \OperatorTok{=}\NormalTok{ node}

    \KeywordTok{def}\NormalTok{ get(}\VariableTok{self}\NormalTok{, key):}
        \ControlFlowTok{if}\NormalTok{ key }\KeywordTok{not} \KeywordTok{in} \VariableTok{self}\NormalTok{.cache: }\ControlFlowTok{return} \OperatorTok{{-}}\DecValTok{1}
\NormalTok{        node }\OperatorTok{=} \VariableTok{self}\NormalTok{.cache[key]}
        \VariableTok{self}\NormalTok{.\_remove(node)}
        \VariableTok{self}\NormalTok{.\_insert\_front(node)}
        \ControlFlowTok{return}\NormalTok{ node.val}

    \KeywordTok{def}\NormalTok{ put(}\VariableTok{self}\NormalTok{, key, val):}
        \ControlFlowTok{if}\NormalTok{ key }\KeywordTok{in} \VariableTok{self}\NormalTok{.cache:}
            \VariableTok{self}\NormalTok{.\_remove(}\VariableTok{self}\NormalTok{.cache[key])}
\NormalTok{        node }\OperatorTok{=}\NormalTok{ DLLNode(key, val)}
        \VariableTok{self}\NormalTok{.cache[key] }\OperatorTok{=}\NormalTok{ node}
        \VariableTok{self}\NormalTok{.\_insert\_front(node)}
        \ControlFlowTok{if} \BuiltInTok{len}\NormalTok{(}\VariableTok{self}\NormalTok{.cache) }\OperatorTok{\textgreater{}} \VariableTok{self}\NormalTok{.cap:}
\NormalTok{            lru }\OperatorTok{=} \VariableTok{self}\NormalTok{.tail.prev}
            \VariableTok{self}\NormalTok{.\_remove(lru)}
            \KeywordTok{del} \VariableTok{self}\NormalTok{.cache[lru.key]}
\end{Highlighting}
\end{Shaded}

\textbf{Interview takeaway:} \textbf{LRU = HashMap + DLL with dummy head and tail. The dummy nodes eliminate all edge-case null checks. LFU is significantly harder --- use HashMap(key\(\rightarrow\)freq) + HashMap(freq\(\rightarrow\)OrderedDict) + min\_freq variable.}

\textbf{Real interview use-case:} LRU cache (literal LeetCode 146), LFU cache (LeetCode 460), design Twitter, design browser history.

\subsection{Segment Tree \& Fenwick Tree (BIT)}\label{segment-tree--fenwick-tree-bit}

\textbf{What it is:}

\begin{itemize}
\tightlist
\item
  \textbf{Segment Tree}: Tree where each node stores aggregate (sum/min/max) over a range. Supports range queries and point/range updates in O(log n).
\item
  \textbf{Fenwick Tree (Binary Indexed Tree)}: Array-based structure using binary representation for prefix sums. Simpler to implement; less general.
\end{itemize}

\textbf{Internal mechanics:}

\setcodelang{}

\begin{verbatim}
Array: [1, 3, 5, 7, 9, 11]  (0-indexed)

Segment Tree (sum):
                  [36: 0-5]
                 /           \
        [9: 0-2]              [27: 3-5]
        /       \             /        \
   [4: 0-1]  [5: 2-2]  [16: 3-4]  [11: 5-5]
    /    \               /     \
[1:0-0][3:1-1]      [7:3-3][9:4-4]

Query sum(1,4): decompose into [1-2] and [3-4] → 8 + 16 = 24
Update arr[2] = 10: update leaf, propagate to root → O(log n)
Space: O(4n) for array representation

Fenwick Tree (BIT) for prefix sums:
  Index:   1  2  3  4  5  6
  arr:    [1, 3, 5, 7, 9,11]
  BIT:    [1, 4, 5,16, 9,20]  ← each stores sum of specific range

  BIT[i] = sum of arr[i - lowbit(i) + 1 ... i]
  lowbit(i) = i & (-i)  ← isolates least significant bit

  Prefix sum query(5):
    i=5 (101): add BIT[5]=9, i=5-(5&-5)=4
    i=4 (100): add BIT[4]=16, i=4-(4&-4)=0 → stop
    Total: 9+16=25 ✓ (but wait: 1+3+5+7+9=25 ✓)

  Update(3, delta=2): propagate up
    i=3: BIT[3]+=2, i=3+(3&-3)=4
    i=4: BIT[4]+=2, i=4+(4&-4)=8 > n → stop
\end{verbatim}

\textbf{Complexity:}

\begin{longtable}[]{@{}
  >{\raggedright\arraybackslash}p{(\columnwidth - 6\tabcolsep) * \real{0.2297}}
  >{\raggedright\arraybackslash}p{(\columnwidth - 6\tabcolsep) * \real{0.2421}}
  >{\raggedright\arraybackslash}p{(\columnwidth - 6\tabcolsep) * \real{0.2641}}
  >{\raggedright\arraybackslash}p{(\columnwidth - 6\tabcolsep) * \real{0.2641}}@{}}
\toprule\noalign{}
\rowcolor{acc!16}
\begin{minipage}[b]{\linewidth}\raggedright
Op
\end{minipage} & \begin{minipage}[b]{\linewidth}\raggedright
Naive Array
\end{minipage} & \begin{minipage}[b]{\linewidth}\raggedright
Segment Tree
\end{minipage} & \begin{minipage}[b]{\linewidth}\raggedright
Fenwick Tree
\end{minipage} \\
\midrule\noalign{}
\endhead
\bottomrule\noalign{}
\endlastfoot
Build & O(n) & O(n) & O(n log n) \\
\rowcolor{zebra}
Point update & O(1) & O(log n) & O(log n) \\
Range update & O(n) & O(log n) & O(log n) \\
\rowcolor{zebra}
Range query & O(n) & O(log n) & O(log n) \\
Space & O(n) & O(4n) & O(n) \\
\end{longtable}

\textbf{When to use:}

\begin{itemize}
\tightlist
\item
  Segment tree: range min/max queries, lazy propagation for range updates
\item
  Fenwick tree: only prefix sums needed; cleaner code
\item
  Both: dynamic range queries where array is mutated
\end{itemize}

\textbf{Key Python patterns:}

\setcodelang{Python}

\begin{Shaded}
\begin{Highlighting}[]
\CommentTok{\# Fenwick Tree (BIT) — prefix sum}
\KeywordTok{class}\NormalTok{ BIT:}
    \KeywordTok{def} \FunctionTok{\_\_init\_\_}\NormalTok{(}\VariableTok{self}\NormalTok{, n):}
        \VariableTok{self}\NormalTok{.n }\OperatorTok{=}\NormalTok{ n}
        \VariableTok{self}\NormalTok{.tree }\OperatorTok{=}\NormalTok{ [}\DecValTok{0}\NormalTok{] }\OperatorTok{*}\NormalTok{ (n }\OperatorTok{+} \DecValTok{1}\NormalTok{)  }\CommentTok{\# 1{-}indexed}

    \KeywordTok{def}\NormalTok{ update(}\VariableTok{self}\NormalTok{, i, delta):    }\CommentTok{\# add delta to position i}
        \ControlFlowTok{while}\NormalTok{ i }\OperatorTok{\textless{}=} \VariableTok{self}\NormalTok{.n:}
            \VariableTok{self}\NormalTok{.tree[i] }\OperatorTok{+=}\NormalTok{ delta}
\NormalTok{            i }\OperatorTok{+=}\NormalTok{ i }\OperatorTok{\&}\NormalTok{ (}\OperatorTok{{-}}\NormalTok{i)          }\CommentTok{\# move to parent}

    \KeywordTok{def}\NormalTok{ query(}\VariableTok{self}\NormalTok{, i):            }\CommentTok{\# prefix sum [1..i]}
\NormalTok{        s }\OperatorTok{=} \DecValTok{0}
        \ControlFlowTok{while}\NormalTok{ i }\OperatorTok{\textgreater{}} \DecValTok{0}\NormalTok{:}
\NormalTok{            s }\OperatorTok{+=} \VariableTok{self}\NormalTok{.tree[i]}
\NormalTok{            i }\OperatorTok{{-}=}\NormalTok{ i }\OperatorTok{\&}\NormalTok{ (}\OperatorTok{{-}}\NormalTok{i)          }\CommentTok{\# move to ancestor}
        \ControlFlowTok{return}\NormalTok{ s}

    \KeywordTok{def}\NormalTok{ range\_query(}\VariableTok{self}\NormalTok{, l, r):   }\CommentTok{\# sum [l..r]}
        \ControlFlowTok{return} \VariableTok{self}\NormalTok{.query(r) }\OperatorTok{{-}} \VariableTok{self}\NormalTok{.query(l }\OperatorTok{{-}} \DecValTok{1}\NormalTok{)}

\CommentTok{\# Segment Tree (iterative, cleaner)}
\KeywordTok{class}\NormalTok{ SegmentTree:}
    \KeywordTok{def} \FunctionTok{\_\_init\_\_}\NormalTok{(}\VariableTok{self}\NormalTok{, nums):}
        \VariableTok{self}\NormalTok{.n }\OperatorTok{=} \BuiltInTok{len}\NormalTok{(nums)}
        \VariableTok{self}\NormalTok{.tree }\OperatorTok{=}\NormalTok{ [}\DecValTok{0}\NormalTok{] }\OperatorTok{*}\NormalTok{ (}\DecValTok{2} \OperatorTok{*} \VariableTok{self}\NormalTok{.n)}
        \CommentTok{\# Build}
        \ControlFlowTok{for}\NormalTok{ i, v }\KeywordTok{in} \BuiltInTok{enumerate}\NormalTok{(nums):}
            \VariableTok{self}\NormalTok{.tree[}\VariableTok{self}\NormalTok{.n }\OperatorTok{+}\NormalTok{ i] }\OperatorTok{=}\NormalTok{ v}
        \ControlFlowTok{for}\NormalTok{ i }\KeywordTok{in} \BuiltInTok{range}\NormalTok{(}\VariableTok{self}\NormalTok{.n }\OperatorTok{{-}} \DecValTok{1}\NormalTok{, }\DecValTok{0}\NormalTok{, }\OperatorTok{{-}}\DecValTok{1}\NormalTok{):}
            \VariableTok{self}\NormalTok{.tree[i] }\OperatorTok{=} \VariableTok{self}\NormalTok{.tree[}\DecValTok{2}\OperatorTok{*}\NormalTok{i] }\OperatorTok{+} \VariableTok{self}\NormalTok{.tree[}\DecValTok{2}\OperatorTok{*}\NormalTok{i}\OperatorTok{+}\DecValTok{1}\NormalTok{]}

    \KeywordTok{def}\NormalTok{ update(}\VariableTok{self}\NormalTok{, i, val):}
\NormalTok{        i }\OperatorTok{+=} \VariableTok{self}\NormalTok{.n}
        \VariableTok{self}\NormalTok{.tree[i] }\OperatorTok{=}\NormalTok{ val}
        \ControlFlowTok{while}\NormalTok{ i }\OperatorTok{\textgreater{}} \DecValTok{1}\NormalTok{:}
\NormalTok{            i }\OperatorTok{//=} \DecValTok{2}
            \VariableTok{self}\NormalTok{.tree[i] }\OperatorTok{=} \VariableTok{self}\NormalTok{.tree[}\DecValTok{2}\OperatorTok{*}\NormalTok{i] }\OperatorTok{+} \VariableTok{self}\NormalTok{.tree[}\DecValTok{2}\OperatorTok{*}\NormalTok{i}\OperatorTok{+}\DecValTok{1}\NormalTok{]}

    \KeywordTok{def}\NormalTok{ query(}\VariableTok{self}\NormalTok{, l, r):         }\CommentTok{\# sum [l, r) half{-}open}
\NormalTok{        res }\OperatorTok{=} \DecValTok{0}
\NormalTok{        l }\OperatorTok{+=} \VariableTok{self}\NormalTok{.n}\OperatorTok{;}\NormalTok{ r }\OperatorTok{+=} \VariableTok{self}\NormalTok{.n}
        \ControlFlowTok{while}\NormalTok{ l }\OperatorTok{\textless{}}\NormalTok{ r:}
            \ControlFlowTok{if}\NormalTok{ l }\OperatorTok{\&} \DecValTok{1}\NormalTok{: res }\OperatorTok{+=} \VariableTok{self}\NormalTok{.tree[l]}\OperatorTok{;}\NormalTok{ l }\OperatorTok{+=} \DecValTok{1}
            \ControlFlowTok{if}\NormalTok{ r }\OperatorTok{\&} \DecValTok{1}\NormalTok{: r }\OperatorTok{{-}=} \DecValTok{1}\OperatorTok{;}\NormalTok{ res }\OperatorTok{+=} \VariableTok{self}\NormalTok{.tree[r]}
\NormalTok{            l }\OperatorTok{//=} \DecValTok{2}\OperatorTok{;}\NormalTok{ r }\OperatorTok{//=} \DecValTok{2}
        \ControlFlowTok{return}\NormalTok{ res}
\end{Highlighting}
\end{Shaded}

\textbf{Interview takeaway:} \textbf{Fenwick tree is simpler to code and sufficient for prefix sum + point update problems. Use segment tree when you need range updates or range min/max. The "count of smaller numbers after self" problem \(\rightarrow\) BIT is the elegant solution.}

\textbf{Real interview use-case:} Range sum query (mutable), count of smaller numbers after self, reverse pairs, falling squares.

\subsection{Bloom Filter}\label{bloom-filter}

\textbf{What it is:} A probabilistic, space-efficient data structure that answers "is X possibly in the set?" or "is X definitely NOT in the set?" --- with false positives but no false negatives.

\textbf{Internal mechanics:}

\setcodelang{}

\begin{verbatim}
Bloom Filter (m=20 bits, k=3 hash functions):

  Initial state:
  [0][0][0][0][0][0][0][0][0][0][0][0][0][0][0][0][0][0][0][0]

  Insert "cat" (hash1=3, hash2=7, hash3=15):
  [0][0][0][1][0][0][0][1][0][0][0][0][0][0][0][1][0][0][0][0]
                  ↑               ↑               ↑

  Insert "dog" (hash1=3, hash2=10, hash3=17):
  [0][0][0][1][0][0][0][1][0][0][1][0][0][0][0][1][0][1][0][0]
  Note: bit 3 now set by both "cat" and "dog"!

  Query "cat": check bits 3,7,15 → all 1 → POSSIBLY in set ✓
  Query "hat": hash1=3, hash2=10, hash3=17 → all 1 → POSSIBLY in set ✗ FALSE POSITIVE!
  Query "xyz": hash1=1 → 0 → DEFINITELY NOT in set ✓

  False positive probability ≈ (1 - e^(-kn/m))^k
  where n=items inserted, m=bit array size, k=hash functions

  Key: no false negatives (if bloom says NO, it's definitely NO)
       cannot delete items (counter-based bloom filters can)
       cannot enumerate stored elements
\end{verbatim}

\textbf{Complexity:}

\begin{longtable}[]{@{}
  >{\raggedright\arraybackslash}p{(\columnwidth - 4\tabcolsep) * \real{0.3670}}
  >{\raggedright\arraybackslash}p{(\columnwidth - 4\tabcolsep) * \real{0.2813}}
  >{\raggedright\arraybackslash}p{(\columnwidth - 4\tabcolsep) * \real{0.3517}}@{}}
\toprule\noalign{}
\rowcolor{acc!16}
\begin{minipage}[b]{\linewidth}\raggedright
Op
\end{minipage} & \begin{minipage}[b]{\linewidth}\raggedright
Time
\end{minipage} & \begin{minipage}[b]{\linewidth}\raggedright
Space
\end{minipage} \\
\midrule\noalign{}
\endhead
\bottomrule\noalign{}
\endlastfoot
Insert & O(k) & O(m) \\
\rowcolor{zebra}
Query & O(k) & O(1) \\
Delete & N/A & - \\
\end{longtable}

\textbf{When to use:}

\begin{itemize}
\tightlist
\item
  Cache pre-filtering: "has this URL been seen before?" before hitting DB
\item
  Distributed systems: "does key exist on this node?" before network call
\item
  Spell checkers: fast "not a word" filter
\item
  HBase/Cassandra: skip SSTable reads for non-existent keys
\end{itemize}

\textbf{Ties to systems:}

\begin{itemize}
\tightlist
\item
  Cassandra uses bloom filters to avoid unnecessary disk I/O
\item
  Google Chrome used bloom filter for Safe Browsing URL blacklist
\item
  Bitcoin uses bloom filters in SPV (Simplified Payment Verification)
\end{itemize}

\textbf{Interview takeaway:} \textbf{Bloom filters trade accuracy for space and speed. False positives OK; false negatives are impossible. When an interviewer asks "how would you check if a URL was seen before using O(1) space?" --- bloom filter is the answer. Note: you can NOT delete from a standard bloom filter.}

\section{Choosing the Right Data Structure}\label{choosing-the-right-data-structure}

Decision table: "I need X \(\rightarrow\) use Y"

\begin{longtable}[]{@{}
  >{\raggedright\arraybackslash}p{(\columnwidth - 4\tabcolsep) * \real{0.3696}}
  >{\raggedright\arraybackslash}p{(\columnwidth - 4\tabcolsep) * \real{0.2935}}
  >{\raggedright\arraybackslash}p{(\columnwidth - 4\tabcolsep) * \real{0.3370}}@{}}
\toprule\noalign{}
\rowcolor{acc!16}
\begin{minipage}[b]{\linewidth}\raggedright
Need
\end{minipage} & \begin{minipage}[b]{\linewidth}\raggedright
Use
\end{minipage} & \begin{minipage}[b]{\linewidth}\raggedright
Why
\end{minipage} \\
\midrule\noalign{}
\endhead
\bottomrule\noalign{}
\endlastfoot
O(1) access by index & Array / List & Contiguous memory \\
\rowcolor{zebra}
O(1) key-value lookup & Hash Map & Hash function \\
Ordered key-value, range queries & BST / Sorted Map & In-order traversal \\
\rowcolor{zebra}
Get min/max repeatedly & Heap & Root always min/max \\
Get kth smallest/largest & Heap of size k & O(n log k) \\
\rowcolor{zebra}
Prefix search / autocomplete & Trie & Shared prefix compression \\
LIFO (undo, recursion) & Stack & Push/pop from one end \\
\rowcolor{zebra}
FIFO (BFS, scheduling) & Queue / Deque & Enqueue/dequeue O(1) \\
Sliding window max/min & Monotonic Deque & Maintain monotone order \\
\rowcolor{zebra}
Next greater/smaller element & Monotonic Stack & Classic pattern \\
Connected components & Union-Find & Near-O(1) union/find \\
\rowcolor{zebra}
Shortest path (unweighted) & BFS (Queue) & Level-by-level expansion \\
Shortest path (weighted, non-neg) & Dijkstra (Heap) & Greedy + priority queue \\
\rowcolor{zebra}
Shortest path (negative weights) & Bellman-Ford & DP approach \\
All pairs shortest path & Floyd-Warshall & O(V\(^3\)) DP \\
\rowcolor{zebra}
Range sum query + updates & Fenwick Tree / Segment Tree & O(log n) both \\
Range min/max query + updates & Segment Tree & BIT can\textquotesingle t do min/max \\
\rowcolor{zebra}
O(1) LRU eviction & HashMap + DLL & O(1) access + O(1) order \\
Membership test, space-constrained & Bloom Filter & Probabilistic, tiny space \\
\rowcolor{zebra}
Count frequencies & HashMap / Counter & Hash map with integer values \\
Find duplicates & HashSet or sorting & O(n) with hash, O(n log n) sort \\
\rowcolor{zebra}
Merge k sorted sequences & Min-Heap of k elements & O(n log k) \\
Detect cycle (undirected graph) & Union-Find & Simpler than DFS coloring \\
\rowcolor{zebra}
Detect cycle (directed graph) & DFS with in-stack tracking & Union-Find doesn\textquotesingle t work here \\
Topological ordering & Kahn\textquotesingle s (BFS) or DFS & Works on DAGs only \\
\end{longtable}

\section{Language Notes}\label{language-notes}

\subsection{Python}\label{python}

\begin{longtable}[]{@{}
  >{\raggedright\arraybackslash}p{(\columnwidth - 4\tabcolsep) * \real{0.2821}}
  >{\raggedright\arraybackslash}p{(\columnwidth - 4\tabcolsep) * \real{0.3462}}
  >{\raggedright\arraybackslash}p{(\columnwidth - 4\tabcolsep) * \real{0.3718}}@{}}
\toprule\noalign{}
\rowcolor{acc!16}
\begin{minipage}[b]{\linewidth}\raggedright
Need
\end{minipage} & \begin{minipage}[b]{\linewidth}\raggedright
Python
\end{minipage} & \begin{minipage}[b]{\linewidth}\raggedright
Notes
\end{minipage} \\
\midrule\noalign{}
\endhead
\bottomrule\noalign{}
\endlastfoot
Dynamic array & \texttt{list} & \texttt{append} O(1) amortized \\
\rowcolor{zebra}
Hash map & \texttt{dict} & \texttt{defaultdict(list/int/set)} \\
Hash set & \texttt{set} & \texttt{add}, \texttt{in} O(1) \\
\rowcolor{zebra}
Ordered map (by key) & \texttt{sortedcontainers.SortedList} & Not stdlib; use \texttt{bisect} + list \\
Min heap & \texttt{heapq} (module) & Negate for max heap \\
\rowcolor{zebra}
Queue / Deque & \texttt{collections.deque} & \texttt{appendleft}, \texttt{popleft} O(1) \\
Counter & \texttt{collections.Counter} & \texttt{most\_common(k)} \\
\rowcolor{zebra}
LRU built-in & \texttt{collections.OrderedDict} & \texttt{move\_to\_end}, \texttt{popitem} \\
Bisect (sorted insert) & \texttt{bisect.bisect\_left/right} & Binary search on sorted list \\
\rowcolor{zebra}
Infinity & \texttt{float(\textquotesingle{}inf\textquotesingle{})} & \\
Char to int & \texttt{ord(\textquotesingle{}a\textquotesingle{})} \(\rightarrow\) 97 & \texttt{chr(97)} \(\rightarrow\) \textquotesingle a\textquotesingle{} \\
\end{longtable}

\setcodelang{Python}

\begin{Shaded}
\begin{Highlighting}[]
\CommentTok{\# Common Python idioms}
\ImportTok{from}\NormalTok{ collections }\ImportTok{import}\NormalTok{ defaultdict, Counter, deque, OrderedDict}
\ImportTok{import}\NormalTok{ heapq}
\ImportTok{from}\NormalTok{ bisect }\ImportTok{import}\NormalTok{ bisect\_left, bisect\_right, insort}

\CommentTok{\# Heap with custom key (push tuples)}
\NormalTok{heapq.heappush(heap, (priority, item))}

\CommentTok{\# Sorted list simulation with bisect}
\ImportTok{import}\NormalTok{ bisect}
\NormalTok{sorted\_list }\OperatorTok{=}\NormalTok{ []}
\NormalTok{bisect.insort(sorted\_list, }\DecValTok{5}\NormalTok{)         }\CommentTok{\# insert and maintain order}
\NormalTok{idx }\OperatorTok{=}\NormalTok{ bisect.bisect\_left(sorted\_list, }\DecValTok{5}\NormalTok{)  }\CommentTok{\# binary search}

\CommentTok{\# defaultdict prevents KeyError}
\NormalTok{graph }\OperatorTok{=}\NormalTok{ defaultdict(}\BuiltInTok{list}\NormalTok{)}
\NormalTok{count }\OperatorTok{=}\NormalTok{ defaultdict(}\BuiltInTok{int}\NormalTok{)}
\end{Highlighting}
\end{Shaded}

\subsection{Java}\label{java}

\begin{longtable}[]{@{}
  >{\raggedright\arraybackslash}p{(\columnwidth - 4\tabcolsep) * \real{0.2372}}
  >{\raggedright\arraybackslash}p{(\columnwidth - 4\tabcolsep) * \real{0.2135}}
  >{\raggedright\arraybackslash}p{(\columnwidth - 4\tabcolsep) * \real{0.5492}}@{}}
\toprule\noalign{}
\rowcolor{acc!16}
\begin{minipage}[b]{\linewidth}\raggedright
Need
\end{minipage} & \begin{minipage}[b]{\linewidth}\raggedright
Java Class
\end{minipage} & \begin{minipage}[b]{\linewidth}\raggedright
Notes
\end{minipage} \\
\midrule\noalign{}
\endhead
\bottomrule\noalign{}
\endlastfoot
Dynamic array & \texttt{ArrayList\textless{}T\textgreater{}} & \texttt{add()} O(1) amortized \\
\rowcolor{zebra}
Hash map & \texttt{HashMap\textless{}K,V\textgreater{}} & Not ordered \\
Ordered map (by key) & \texttt{TreeMap\textless{}K,V\textgreater{}} & Red-Black tree, O(log n) ops \\
\rowcolor{zebra}
Ordered set & \texttt{TreeSet\textless{}T\textgreater{}} & Red-Black tree \\
Hash set & \texttt{HashSet\textless{}T\textgreater{}} & \\
\rowcolor{zebra}
Linked hash map & \texttt{LinkedHashMap\textless{}K,V\textgreater{}} & Insertion-order iteration (LRU) \\
Min/Max heap & \texttt{PriorityQueue\textless{}T\textgreater{}} & Min by default; pass \texttt{Comparator.reverseOrder()} for max \\
\rowcolor{zebra}
Stack & \texttt{Deque\textless{}T\textgreater{}} as stack & Prefer \texttt{ArrayDeque} over \texttt{Stack} \\
Queue & \texttt{ArrayDeque\textless{}T\textgreater{}} & Or \texttt{LinkedList\textless{}T\textgreater{}} \\
\rowcolor{zebra}
Deque & \texttt{ArrayDeque\textless{}T\textgreater{}} & \\
\end{longtable}

\setcodelang{Java}

\begin{Shaded}
\begin{Highlighting}[]
\CommentTok{// Java idioms}
\BuiltInTok{PriorityQueue}\OperatorTok{\textless{}}\BuiltInTok{Integer}\OperatorTok{\textgreater{}}\NormalTok{ minHeap }\OperatorTok{=} \KeywordTok{new} \BuiltInTok{PriorityQueue}\OperatorTok{\textless{}\textgreater{}();}
\BuiltInTok{PriorityQueue}\OperatorTok{\textless{}}\BuiltInTok{Integer}\OperatorTok{\textgreater{}}\NormalTok{ maxHeap }\OperatorTok{=} \KeywordTok{new} \BuiltInTok{PriorityQueue}\OperatorTok{\textless{}\textgreater{}(}\BuiltInTok{Collections}\OperatorTok{.}\FunctionTok{reverseOrder}\OperatorTok{());}

\BuiltInTok{Map}\OperatorTok{\textless{}}\BuiltInTok{String}\OperatorTok{,} \BuiltInTok{Integer}\OperatorTok{\textgreater{}}\NormalTok{ map }\OperatorTok{=} \KeywordTok{new} \BuiltInTok{HashMap}\OperatorTok{\textless{}\textgreater{}();}
\NormalTok{map}\OperatorTok{.}\FunctionTok{getOrDefault}\OperatorTok{(}\NormalTok{key}\OperatorTok{,} \DecValTok{0}\OperatorTok{);}
\NormalTok{map}\OperatorTok{.}\FunctionTok{computeIfAbsent}\OperatorTok{(}\NormalTok{key}\OperatorTok{,}\NormalTok{ k }\OperatorTok{{-}\textgreater{}} \KeywordTok{new} \BuiltInTok{ArrayList}\OperatorTok{\textless{}\textgreater{}()).}\FunctionTok{add}\OperatorTok{(}\NormalTok{val}\OperatorTok{);}

\BuiltInTok{Deque}\OperatorTok{\textless{}}\BuiltInTok{Integer}\OperatorTok{\textgreater{}}\NormalTok{ stack }\OperatorTok{=} \KeywordTok{new} \BuiltInTok{ArrayDeque}\OperatorTok{\textless{}\textgreater{}();}  \CommentTok{// use as stack}
\NormalTok{stack}\OperatorTok{.}\FunctionTok{push}\OperatorTok{(}\DecValTok{1}\OperatorTok{);}\NormalTok{ stack}\OperatorTok{.}\FunctionTok{pop}\OperatorTok{();}\NormalTok{ stack}\OperatorTok{.}\FunctionTok{peek}\OperatorTok{();}

\BuiltInTok{Queue}\OperatorTok{\textless{}}\BuiltInTok{Integer}\OperatorTok{\textgreater{}}\NormalTok{ queue }\OperatorTok{=} \KeywordTok{new} \BuiltInTok{ArrayDeque}\OperatorTok{\textless{}\textgreater{}();}  \CommentTok{// use as queue}
\NormalTok{queue}\OperatorTok{.}\FunctionTok{offer}\OperatorTok{(}\DecValTok{1}\OperatorTok{);}\NormalTok{ queue}\OperatorTok{.}\FunctionTok{poll}\OperatorTok{();}\NormalTok{ queue}\OperatorTok{.}\FunctionTok{peek}\OperatorTok{();}

\BuiltInTok{TreeMap}\OperatorTok{\textless{}}\BuiltInTok{Integer}\OperatorTok{,} \BuiltInTok{Integer}\OperatorTok{\textgreater{}}\NormalTok{ ordered }\OperatorTok{=} \KeywordTok{new} \BuiltInTok{TreeMap}\OperatorTok{\textless{}\textgreater{}();}
\NormalTok{ordered}\OperatorTok{.}\FunctionTok{floorKey}\OperatorTok{(}\NormalTok{x}\OperatorTok{);}    \CommentTok{// largest key ≤ x}
\NormalTok{ordered}\OperatorTok{.}\FunctionTok{ceilingKey}\OperatorTok{(}\NormalTok{x}\OperatorTok{);}  \CommentTok{// smallest key ≥ x}
\end{Highlighting}
\end{Shaded}

\subsection{C++}\label{c}

\begin{longtable}[]{@{}
  >{\raggedright\arraybackslash}p{(\columnwidth - 4\tabcolsep) * \real{0.2381}}
  >{\raggedright\arraybackslash}p{(\columnwidth - 4\tabcolsep) * \real{0.4524}}
  >{\raggedright\arraybackslash}p{(\columnwidth - 4\tabcolsep) * \real{0.3095}}@{}}
\toprule\noalign{}
\rowcolor{acc!16}
\begin{minipage}[b]{\linewidth}\raggedright
Need
\end{minipage} & \begin{minipage}[b]{\linewidth}\raggedright
C++ Container
\end{minipage} & \begin{minipage}[b]{\linewidth}\raggedright
Notes
\end{minipage} \\
\midrule\noalign{}
\endhead
\bottomrule\noalign{}
\endlastfoot
Dynamic array & \texttt{vector\textless{}T\textgreater{}} & \texttt{push\_back} O(1) amortized \\
\rowcolor{zebra}
Hash map & \texttt{unordered\_map\textless{}K,V\textgreater{}} & O(1) avg \\
Ordered map (by key) & \texttt{map\textless{}K,V\textgreater{}} & Red-Black tree, O(log n) \\
\rowcolor{zebra}
Hash set & \texttt{unordered\_set\textless{}T\textgreater{}} & \\
Ordered set & \texttt{set\textless{}T\textgreater{}} & + \texttt{lower\_bound}, \texttt{upper\_bound} \\
\rowcolor{zebra}
Priority queue (max) & \texttt{priority\_queue\textless{}T\textgreater{}} & Max by default \\
Priority queue (min) & \texttt{priority\_queue\textless{}T,vector\textless{}T\textgreater{},greater\textless{}T\textgreater{}\textgreater{}} & \\
\rowcolor{zebra}
Stack & \texttt{stack\textless{}T\textgreater{}} & \\
Queue & \texttt{queue\textless{}T\textgreater{}} & \\
\rowcolor{zebra}
Deque & \texttt{deque\textless{}T\textgreater{}} & \\
\end{longtable}

\setcodelang{C++}

\begin{Shaded}
\begin{Highlighting}[]
\CommentTok{// C++ idioms}
\PreprocessorTok{\#include }\ImportTok{\textless{}bits/stdc++.h\textgreater{}}
\KeywordTok{using} \KeywordTok{namespace}\NormalTok{ std}\OperatorTok{;}

\NormalTok{unordered\_map}\OperatorTok{\textless{}}\DataTypeTok{int}\OperatorTok{,}\DataTypeTok{int}\OperatorTok{\textgreater{}}\NormalTok{ freq}\OperatorTok{;}
\NormalTok{freq}\OperatorTok{[}\NormalTok{x}\OperatorTok{]++;}  \CommentTok{// default{-}initializes to 0}

\CommentTok{// Min heap}
\NormalTok{priority\_queue}\OperatorTok{\textless{}}\DataTypeTok{int}\OperatorTok{,}\NormalTok{ vector}\OperatorTok{\textless{}}\DataTypeTok{int}\OperatorTok{\textgreater{},}\NormalTok{ greater}\OperatorTok{\textless{}}\DataTypeTok{int}\OperatorTok{\textgreater{}\textgreater{}}\NormalTok{ minHeap}\OperatorTok{;}

\CommentTok{// Ordered set with lower\_bound}
\NormalTok{set}\OperatorTok{\textless{}}\DataTypeTok{int}\OperatorTok{\textgreater{}}\NormalTok{ s}\OperatorTok{;}
\KeywordTok{auto}\NormalTok{ it }\OperatorTok{=}\NormalTok{ s}\OperatorTok{.}\NormalTok{lower\_bound}\OperatorTok{(}\NormalTok{x}\OperatorTok{);}  \CommentTok{// O(log n)}

\CommentTok{// Sort with custom comparator}
\NormalTok{sort}\OperatorTok{(}\NormalTok{v}\OperatorTok{.}\NormalTok{begin}\OperatorTok{(),}\NormalTok{ v}\OperatorTok{.}\NormalTok{end}\OperatorTok{(),} \OperatorTok{[](}\KeywordTok{auto}\OperatorTok{\&}\NormalTok{ a}\OperatorTok{,} \KeywordTok{auto}\OperatorTok{\&}\NormalTok{ b}\OperatorTok{)\{} \ControlFlowTok{return}\NormalTok{ a}\OperatorTok{[}\DecValTok{0}\OperatorTok{]} \OperatorTok{\textless{}}\NormalTok{ b}\OperatorTok{[}\DecValTok{0}\OperatorTok{];} \OperatorTok{\});}
\end{Highlighting}
\end{Shaded}

\section{Real-Life Analogies}\label{real-life-analogies}

\begin{longtable}[]{@{}
  >{\raggedright\arraybackslash}p{(\columnwidth - 2\tabcolsep) * \real{0.2313}}
  >{\raggedright\arraybackslash}p{(\columnwidth - 2\tabcolsep) * \real{0.7687}}@{}}
\toprule\noalign{}
\rowcolor{acc!16}
\begin{minipage}[b]{\linewidth}\raggedright
Data Structure
\end{minipage} & \begin{minipage}[b]{\linewidth}\raggedright
Analogy
\end{minipage} \\
\midrule\noalign{}
\endhead
\bottomrule\noalign{}
\endlastfoot
Array & Numbered mailboxes in a row --- access box 42 instantly \\
\rowcolor{zebra}
Hash Map & Library card catalog --- look up by title, get shelf location \\
Linked List & Train cars --- each car knows only the next car \\
\rowcolor{zebra}
Stack & Stack of plates --- only access the top \\
Queue & Queue at a bank --- first in, first served \\
\rowcolor{zebra}
Deque & Train that can board/exit from both ends \\
BST & Phone book --- organized so you can binary search \\
\rowcolor{zebra}
Heap & Priority ER room --- most critical patient always treated first \\
Trie & Autocomplete on your phone --- shared prefix saves space \\
\rowcolor{zebra}
Graph & Road map --- cities are nodes, roads are edges \\
Union-Find & Social groups --- merging friend circles \\
\rowcolor{zebra}
Segment Tree & Population pyramid --- sum over any age range quickly \\
Fenwick Tree & Odometer with clever bit tricks for partial sums \\
\rowcolor{zebra}
LRU Cache & Small whiteboard --- erase what was written longest ago for space \\
Bloom Filter & Bouncer list --- might let someone in wrongly, never turns away regulars \\
\end{longtable}

\section{Common Interview Questions}\label{common-interview-questions}

\subsection{Arrays}\label{arrays}

\begin{itemize}
\tightlist
\item
  Two Sum (LeetCode 1)
\item
  Best Time to Buy and Sell Stock (LC 121)
\item
  Maximum Subarray / Kadane\textquotesingle s (LC 53)
\item
  Container With Most Water (LC 11)
\item
  Trapping Rain Water (LC 42)
\item
  Rotate Array (LC 189)
\item
  Product of Array Except Self (LC 238)
\item
  Merge Intervals (LC 56)
\end{itemize}

\subsection{Strings}\label{strings-1}

\begin{itemize}
\tightlist
\item
  Longest Substring Without Repeating Characters (LC 3)
\item
  Minimum Window Substring (LC 76)
\item
  Group Anagrams (LC 49)
\item
  Valid Palindrome (LC 125)
\item
  Encode and Decode Strings (LC 271)
\item
  Longest Palindromic Substring (LC 5)
\item
  String to Integer (atoi) (LC 8)
\end{itemize}

\subsection{Hash Maps/Sets}\label{hash-mapssets}

\begin{itemize}
\tightlist
\item
  Two Sum (LC 1)
\item
  Top K Frequent Elements (LC 347)
\item
  Longest Consecutive Sequence (LC 128)
\item
  Subarray Sum Equals K (LC 560)
\item
  Valid Sudoku (LC 36)
\item
  Find All Anagrams in a String (LC 438)
\end{itemize}

\subsection{Linked Lists}\label{linked-lists-1}

\begin{itemize}
\tightlist
\item
  Reverse Linked List (LC 206)
\item
  Merge Two Sorted Lists (LC 21)
\item
  Linked List Cycle (LC 141)
\item
  Reorder List (LC 143)
\item
  Remove Nth Node From End (LC 19)
\item
  Merge K Sorted Lists (LC 23)
\end{itemize}

\subsection{Stacks/Queues}\label{stacksqueues}

\begin{itemize}
\tightlist
\item
  Valid Parentheses (LC 20)
\item
  Min Stack (LC 155)
\item
  Daily Temperatures (LC 739)
\item
  Largest Rectangle in Histogram (LC 84)
\item
  Sliding Window Maximum (LC 239)
\item
  Evaluate Reverse Polish Notation (LC 150)
\end{itemize}

\subsection{Trees}\label{trees-1}

\begin{itemize}
\tightlist
\item
  Maximum Depth of Binary Tree (LC 104)
\item
  Invert Binary Tree (LC 226)
\item
  Same Tree (LC 100)
\item
  Binary Tree Level Order Traversal (LC 102)
\item
  Validate Binary Search Tree (LC 98)
\item
  Lowest Common Ancestor of BST (LC 235)
\item
  Kth Smallest Element in BST (LC 230)
\item
  Serialize and Deserialize Binary Tree (LC 297)
\item
  Diameter of Binary Tree (LC 543)
\end{itemize}

\subsection{Heaps}\label{heaps}

\begin{itemize}
\tightlist
\item
  Kth Largest Element in an Array (LC 215)
\item
  K Closest Points to Origin (LC 973)
\item
  Find Median from Data Stream (LC 295)
\item
  Task Scheduler (LC 621)
\item
  Merge K Sorted Lists (LC 23)
\end{itemize}

\subsection{Tries}\label{tries}

\begin{itemize}
\tightlist
\item
  Implement Trie (LC 208)
\item
  Word Search II (LC 212)
\item
  Replace Words (LC 648)
\item
  Design Search Autocomplete System (LC 642)
\end{itemize}

\subsection{Graphs}\label{graphs-1}

\begin{itemize}
\tightlist
\item
  Number of Islands (LC 200)
\item
  Clone Graph (LC 133)
\item
  Pacific Atlantic Water Flow (LC 417)
\item
  Course Schedule (LC 207) --- cycle detection
\item
  Number of Connected Components (LC 323)
\item
  Word Ladder (LC 127) --- BFS shortest path
\item
  Network Delay Time (LC 743) --- Dijkstra
\item
  Cheapest Flights Within K Stops (LC 787) --- modified Dijkstra/Bellman-Ford
\end{itemize}

\subsection{Union-Find}\label{union-find}

\begin{itemize}
\tightlist
\item
  Number of Provinces (LC 547)
\item
  Redundant Connection (LC 684)
\item
  Accounts Merge (LC 721)
\item
  Satisfiability of Equality Equations (LC 990)
\end{itemize}

\subsection{LRU/LFU}\label{lrulfu}

\begin{itemize}
\tightlist
\item
  LRU Cache (LC 146)
\item
  LFU Cache (LC 460)
\end{itemize}

\subsection{Segment Tree / BIT}\label{segment-tree--bit}

\begin{itemize}
\tightlist
\item
  Range Sum Query --- Mutable (LC 307)
\item
  Count of Smaller Numbers After Self (LC 315)
\item
  Reverse Pairs (LC 493)
\end{itemize}

\section{Typical Mistakes Candidates Make}\label{typical-mistakes-candidates-make}

\begin{enumerate}
\def\labelenumi{\arabic{enumi}.}
\item
  \textbf{Using \texttt{list.pop(0)} for queues} --- It\textquotesingle s O(n). Always use \texttt{collections.deque} and \texttt{popleft()}.
\item
  \textbf{String concatenation in loops} --- \texttt{result\ +=\ s} inside a loop is O(n\(^2\)). Use \texttt{"".join(parts)}.
\item
  \textbf{Not handling the empty/single-element edge case} --- For linked lists, trees, arrays. Check \texttt{if\ not\ root}, \texttt{if\ not\ nums}, \texttt{if\ len(nums)\ \textless{}\ 2} early.
\item
  \textbf{BST validation using only local \texttt{left\ \textless{}\ root\ \textless{}\ right}} --- This misses nodes in subtrees. Pass bounds \texttt{(lo,\ hi)} down the recursion.
\item
  \textbf{Modifying a collection while iterating it} --- Causes RuntimeError. Copy or collect indices first.
\item
  \textbf{Forgetting to mark nodes visited before enqueuing (BFS)} --- Leads to exponential revisits. Mark visited when you enqueue, not when you dequeue.
\item
  \textbf{Using a set as a hash map key} --- Sets are unhashable. Use \texttt{frozenset} or \texttt{tuple(sorted(...))}.
\item
  \textbf{Forgetting Python\textquotesingle s \texttt{heapq} is min-heap only} --- For max-heap, push \texttt{-val} and negate on pop.
\item
  \textbf{O(n\(^2\)) nested loop when O(n) hash map exists} --- When you write two nested loops looking for a pair, always ask: "can I precompute what I\textquotesingle m looking for?"
\item
  \textbf{Not using dummy head node for linked lists} --- Head deletion becomes a special case. Dummy head eliminates it.
\item
  \textbf{Recursion without memoization} --- Fibonacci, path counting, etc. naturally revisit subproblems. Add \texttt{@lru\_cache} or a memo dict.
\item
  \textbf{Union-Find without path compression} --- Without it, find is O(n). Always compress.
\item
  \textbf{Assuming \texttt{dict} maintains insertion order in Python 2} --- It doesn\textquotesingle t. Python 3.7+ guarantees it; don\textquotesingle t rely on this in cross-version code.
\item
  \textbf{Off-by-one in sliding window} --- When window is \texttt{{[}l,\ r{]}}, size is \texttt{r\ -\ l\ +\ 1}. Sketch a small example before coding.
\item
  \textbf{Passing arrays by reference and mutating them} --- In Python lists are passed by reference. If you need to preserve the original, copy with \texttt{arr{[}:{]}}.
\end{enumerate}

\section{Revision Cheat Sheet}\label{revision-cheat-sheet}

\subsection{Complexities to Memorize (Absolute Must)}\label{complexities-to-memorize-absolute-must}

\setcodelang{}

\begin{verbatim}
Array access:     O(1)     | Hash map get/put: O(1) avg
Array search:     O(n)     | BST search:       O(log n) balanced
Append to list:   O(1) A   | Heap push/pop:    O(log n)
List insert/del:  O(n)     | Build heap:       O(n)  ← non-obvious!
Trie ops:         O(m)     | Segment tree:     O(log n) query/update
BFS/DFS:          O(V+E)   | Union-Find:       O(α(n)) ≈ O(1)
Dijkstra:         O((V+E)log V)
Heapsort:         O(n log n)
Sort:             O(n log n)
Binary search:    O(log n)
String join:      O(n)     | String concat loop: O(n²) ← DANGER

A = amortized
\end{verbatim}

\subsection{Decision Hooks (6 questions to ask every interview)}\label{decision-hooks-6-questions-to-ask-every-interview}

\begin{enumerate}
\def\labelenumi{\arabic{enumi}.}
\tightlist
\item
  \textbf{"Do I need O(1) lookup?"} \(\rightarrow\) Hash map or hash set.
\item
  \textbf{"Do I need sorted order or range queries?"} \(\rightarrow\) BST, sorted array + binary search, or segment tree.
\item
  \textbf{"Do I need repeated min/max?"} \(\rightarrow\) Heap.
\item
  \textbf{"Is this a traversal/reachability problem?"} \(\rightarrow\) BFS (shortest path) or DFS (connectivity, cycles).
\item
  \textbf{"Am I connecting/grouping things?"} \(\rightarrow\) Union-Find.
\item
  \textbf{"Do I need prefix operations?"} \(\rightarrow\) Trie (strings) or Fenwick tree (numbers).
\end{enumerate}

\subsection{Quick Pattern Mapping}\label{quick-pattern-mapping}

\begin{longtable}[]{@{}
  >{\raggedright\arraybackslash}p{(\columnwidth - 2\tabcolsep) * \real{0.4906}}
  >{\raggedright\arraybackslash}p{(\columnwidth - 2\tabcolsep) * \real{0.5094}}@{}}
\toprule\noalign{}
\rowcolor{acc!16}
\begin{minipage}[b]{\linewidth}\raggedright
Pattern
\end{minipage} & \begin{minipage}[b]{\linewidth}\raggedright
Structure
\end{minipage} \\
\midrule\noalign{}
\endhead
\bottomrule\noalign{}
\endlastfoot
Sliding window & Array + two pointers \\
\rowcolor{zebra}
Sliding window max/min & Monotonic deque \\
Next greater element & Monotonic stack \\
\rowcolor{zebra}
Level-order traversal & Queue (BFS) \\
Path/exhaustive search & Stack (DFS) / recursion \\
\rowcolor{zebra}
Shortest path (unweighted) & BFS \\
Shortest path (weighted) & Dijkstra (min-heap) \\
\rowcolor{zebra}
Connected components & BFS/DFS or Union-Find \\
Top K elements & Heap of size K \\
\rowcolor{zebra}
Frequency counting & Hash map / Counter \\
Deduplication & Hash set \\
\rowcolor{zebra}
Prefix sum & Fenwick tree / prefix array \\
Range sum (mutable) & Fenwick tree \\
\rowcolor{zebra}
Range min/max (mutable) & Segment tree \\
Word prefix search & Trie \\
\rowcolor{zebra}
Order + O(1) eviction & LRU (HashMap + DLL) \\
Large-scale membership & Bloom filter \\
\end{longtable}

\subsection{The 5 Structures You MUST Be Able to Code From Scratch}\label{the-5-structures-you-must-be-able-to-code-from-scratch}

\begin{enumerate}
\def\labelenumi{\arabic{enumi}.}
\tightlist
\item
  \textbf{Hash Map} (open addressing or chaining --- at least explain internals)
\item
  \textbf{Linked List} with insert, delete, reverse
\item
  \textbf{Binary Heap} with push, pop, heapify
\item
  \textbf{Trie} with insert, search, startsWith
\item
  \textbf{LRU Cache} (HashMap + doubly linked list)
\end{enumerate}

\subsection{Memory Aids}\label{memory-aids}

\begin{itemize}
\tightlist
\item
  \textbf{Heap is NOT a BST} --- heap has parent \(\leq\) children (only), BST has left \textless{} parent \textless{} right.
\item
  \textbf{Trie depth = key length, not n} --- n = number of keys.
\item
  \textbf{Union-Find has no delete} --- it\textquotesingle s append-only by design.
\item
  \textbf{BFS finds shortest path; DFS does not} --- in unweighted graphs.
\item
  \textbf{\texttt{heapify} is O(n)} --- this surprises everyone. Proof: most nodes are at the bottom and need to sift down only a little.
\item
  \textbf{Hash map worst case O(n)} --- when all keys hash to same bucket (adversarial input or bad hash). In practice O(1).
\item
  \textbf{AVL is stricter than Red-Black} --- AVL: max height diff 1; RB: height \(\leq\) 2*log(n). RB faster for insert/delete, AVL faster for lookup.
\end{itemize}

\end{document}
